\documentclass[11pt]{article}
\ifdefined\pdfinfoomitdate\pdfinfoomitdate=1\fi
\ifdefined\pdftrailerid\pdftrailerid{}\fi
\ifdefined\pdfsuppressptexinfo\pdfsuppressptexinfo=15\fi
\usepackage[T1]{fontenc}
\usepackage{lmodern}
\usepackage[margin=1in]{geometry}
\usepackage{amsmath,amssymb,amsthm,amscd,mathtools}
\usepackage{microtype}
\usepackage{enumitem}
\usepackage{array,tabularx,booktabs}
\usepackage{xurl}
\usepackage[hidelinks]{hyperref}
\input{glyphtounicode}
\input{glyphtounicode-cmex}
\pdfgentounicode=1
\hypersetup{
 pdftitle={Counterexamples to weak chromatic splitting: sphere kernels and descent exponents},
 pdfauthor={OpenAI},
 pdfsubject={Weak chromatic splitting for the completed sphere},
 pdfkeywords={chromatic homotopy theory, Morava E-theory, descent exponents, beta family}
}
\setlist[enumerate]{label=(\roman*),leftmargin=*}
\newtheorem{theorem}{Theorem}[section]
\newtheorem{proposition}[theorem]{Proposition}
\newtheorem{lemma}[theorem]{Lemma}
\newtheorem{corollary}[theorem]{Corollary}
\theoremstyle{definition}
\newtheorem{definition}[theorem]{Definition}
\theoremstyle{remark}
\newtheorem{remark}[theorem]{Remark}
\newcommand{\Sp}{\mathrm{Sp}_{(p)}}
\newcommand{\F}{\mathbb F}
\newcommand{\Z}{\mathbb Z}
\newcommand{\Q}{\mathbb Q}
\newcommand{\G}{\mathbb G}
\newcommand{\St}{\mathbb S}
\newcommand{\m}{\mathfrak m}
\newcommand{\one}{\mathbf 1}
\DeclareMathOperator{\fib}{fib}
\DeclareMathOperator{\cofib}{cofib}
\DeclareMathOperator{\Tot}{Tot}
\DeclareMathOperator{\im}{im}
\DeclareMathOperator{\Gal}{Gal}
\DeclareMathOperator{\Fun}{Fun}
\DeclareMathOperator{\Tr}{Tr}
\DeclareMathOperator{\cts}{cts}
\DeclareMathOperator{\cd}{cd}
\DeclareMathOperator{\expn}{exp}

\title{Counterexamples to weak chromatic splitting:\\
       sphere kernels and descent exponents}
\author{OpenAI}
\date{September 27, 2026}
\begin{document}
\maketitle
\begin{abstract}
For the derived \(p\)-completion of the sphere, the canonical weak chromatic
splitting map has no homotopy retraction at height \(p\) for every prime
\(p\geq5\), and at height \(p+1\) for every prime \(p\geq7\).
The classical sphere product \(\beta_1^{(p-1)^2}\) gives a nonzero
kernel class in both ranges.
\end{abstract}
\tableofcontents
\section{The canonical weak splitting map}\label{sec:introduction}

Fix a prime $p$ and work in ordinary $p$-local spectra, with sphere
$S=S^0_{(p)}$. Let $K(a)$ denote Morava $K$-theory of height $a$, put
$K(0)=H\Q$, and write
\[
 L_r=L_{K(0)\vee\cdots\vee K(r)},\qquad
 X=S_p^\wedge=\operatorname{holim}_{j\geq1}S/p^j.
\]
Thus $L_r$ is Johnson--Wilson $E(r)$-localization. For a spectrum $Z$,
let $\eta_Z^{(n)}:Z\to L_{K(n)}Z$ be the localization unit. We study
\begin{equation}\label{eq:weak-unit}
 i_{n,X}=L_{n-1}(\eta_X^{(n)}):
 L_{n-1}X\longrightarrow L_{n-1}L_{K(n)}X.
\end{equation}
The weak chromatic splitting conjecture predicts a homotopy retraction
of this map whenever $X$ is the derived $p$-completion of a finite
spectrum. A retraction is a map of spectra $r$ with
$r i_{n,X}\simeq\mathrm{id}$; it need not be natural in the finite input.
We disprove this universal assertion with finite input the sphere.

For an odd prime, let
\[
 \beta_1\in\pi_{d_p}S,\qquad
 d_p=2p(p-1)-2,\qquad k_p=(p-1)^2,
 \qquad z_p=\beta_1^{k_p},
\]
where $\beta_1$ is the classical first beta element.

\begin{theorem}[A named kernel at height $p$]
\label{thm:height-p-kernel}
For every prime $p\geq5$, the image of $z_p$ in
$\pi_{d_pk_p}L_{p-1}X$ is nonzero and is killed by
$\pi_{d_pk_p}(i_{p,X})$. Consequently $i_{p,X}$ has no homotopy
retraction.
\end{theorem}

\begin{theorem}[The same kernel at height $p+1$]
\label{thm:height-plus-one-kernel}
For every prime $p\geq7$, the image of the same $z_p$ in
$\pi_{d_pk_p}L_pX$ is nonzero and is killed by
$\pi_{d_pk_p}(i_{p+1,X})$. Consequently $i_{p+1,X}$ has no homotopy
retraction.
\end{theorem}

\begin{corollary}[The prime-five class]\label{thm:five}
At $p=5$, the image of $\beta_1^{16}$ in $\pi_{608}L_4X$ is nonzero
and belongs to the kernel of $\pi_{608}(i_{5,X})$.
\end{corollary}

The proof detects $z_p$ under a height-$(p-1)$ unit and kills it under
both upper units. The lower detection follows from the Hopkins--Miller
calculation recorded by Heard
\cite[proof of Theorem~5.6 and Lemma~5.8]{heard2015}: the actual unit
detects the named $\beta_1$, and its $k_p$th power remains nonzero in
the ordinary homotopy fixed point spectral sequence. At the upper
height $a$, we construct a separated multiplicative filtration with
the unit image of $\beta_1$ in filtration at least two and
$F^{a^2+2}=0$. The required
inequalities are
\[
 2k_p-(p^2+1)=p^2-4p+1>0\quad(p\geq5),
\]
\[
 2k_p-((p+1)^2+1)=p(p-6)>0\quad(p\geq7).
\]
Completion and naturality then place the nonzero lower image in the
kernel of the exact map \eqref{eq:weak-unit}. For example, at $p=7$
the class $\beta_1^{36}$ has degree $82\cdot36=2952$ and gives
kernels at both heights seven and eight.

\paragraph{Notation.}\label{par:common-notation}
Let $E_a$ be Morava $E$-theory of the height-$a$ Honda formal group
over $\F_{p^a}$, and put $B_a=L_{K(a)}S$. Unless a localized tensor
is explicitly displayed, smash products and function spectra are
ordinary. All homotopy groups are those of the underlying spectra.

\subsection{An independent descent-exponent obstruction}

The height-$p$ nonretraction also follows from incompatible finite
descent lengths. Put $t=p-1$, and let the finite pure-stabilizer subgroup
$H=C_p\rtimes C_{(p-1)^2}$ act on $E_t$. In
$\mathcal D_H=\Fun(BH,\Sp)$, with pointwise ordinary smash, put
$A_H=F(H_+,S)$ with its translation action and
$J_H=\fib(S\to A_H)$. Define $\expn_H(V)$ as the least $b\geq1$
such that $J_H^{\wedge b}\wedge V\to V$ is null, and as infinity
if no such $b$ exists. Set
\[
 T(V)=L_{K(t)}(E_t\wedge V),
\]
with $H$ acting on the first factor.

\begin{theorem}[Incompatible descent exponents]\label{thm:exponents}
For every prime $p\geq5$, with $m=p^2+2$,
\[
 \expn_H(E_{p-1})>m,
 \qquad \expn_H\bigl(T(B_p)\bigr)\leq m.
\]
A homotopy retraction of $i_{p,X}$ would make $E_{p-1}$ an
$H$-equivariant retract of $T(B_p)$. Hence $i_{p,X}$ has no homotopy
retraction, independently of Theorem~\ref{thm:height-p-kernel}.
\end{theorem}

For the lower bound, a class survives to one specified finite page of
the ordinary homotopy fixed point spectral sequence. No permanent
cycle or sphere representative is needed. For the upper bound,
Devinatz--Hopkins' nilpotence theorem supplies a finite exponent before
any homotopy test is chosen \cite[Appendix~A, Proposition~A.3]{dh2004}.
The numerical cochain bound then makes the relevant map zero after
every finite ordinary test. A Brown--Comenetz argument proves that
the map itself is null, and a coherent trace construction transports
the resulting finite filtration to height $p-1$. The finite-stage
formalism used here belongs to the nilpotence and descent theory of
Mathew and Mathew--Naumann--Noel
\cite{mathew2016,mnn2017,mnn2019}.

\subsection{Relation to earlier splitting results}

Hopkins' chromatic splitting conjecture, recorded by Hovey
\cite[Conjecture~4.2, especially part~(v)]{hovey1995}, predicts a
specific decomposition of chromatic overlap as well as the weak unit
retraction. These predictions differ already at $n=p=2$.
Beaudry disproves the proposed strong summand formula
\cite[Theorem~1.0.4]{beaudry2017}; Beaudry--Goerss--Henn give a revised
decomposition with Moore terms that retains the canonical unit as a
retract \cite[Theorem~1.1.6]{bgh2022}. Their
Conjectures~1.1.10--1.1.11 distinguish the strong and
completed-finite-input weak formulations, and their introduction
records the weak form at all primes through height two. The theorem
numbers refer to the versioned texts identified in the bibliography.
Outside finite inputs, Minami reports Devinatz's counterexample for
a $BP$ analogue allowing completed $BP$
\cite[Introduction and Remark~1.1]{minami2003}. The present
counterexamples use the completion of the finite sphere.

Long beta powers in the sphere were known earlier. Lee--Ravenel prove
$\beta_1^{p^2-p-1}\ne0$ for $p\geq7$ and report
$\beta_1^{17}\ne0$, $\beta_1^{18}=0$ at five
\cite[opening theorem and introduction]{lee-ravenel1994}.
Ravenel's first-beta detection calculation relates the cobar class to
stabilizer cohomology \cite[Theorem~4 and its proof in Section~2]{ravenel1978}.
The additional requirement here is detection after a specified lower
localization and vanishing of that same named sphere product under
the upper unit.

\paragraph{Organization.}
Sections~\ref{sec:completion}--\ref{sec:kernels} prove the named
kernel results. Their upper bound uses dimension at height $p+1$
and a semilinear trace-one contraction at height $p$, retaining the
possible inverse-limit degree in both cases.
Sections~\ref{sec:towers}--\ref{sec:exponent-transport} give the
independent exponent proof. Appendix~\ref{sec:ordinary-products}
proves the ordinary filtered-product lemma used by the lower detector
and by the finite-page argument.

\section{Completion and the canonical map}\label{sec:completion}
The lower detectors are local spectra, whereas the weak splitting problem
uses a completed finite input. The following argument connects these
settings without interchanging smash products and inverse limits.

\begin{lemma}\label{lem:completion}
The completion map $c:S\to X=S_p^\wedge$ is a $K(a)$-equivalence for
every integer $a\geq1$. If $1\leq t\leq r$ and $Y$ is $K(t)$-local,
every map $S\to Y$ factors through $S\to X\to L_rX$.
\end{lemma}
\begin{proof}
Take the inverse limit of the cofiber sequences
$S\xrightarrow{p^j}S\to S/p^j$, with multiplication by $p$ on the
left spheres and identity on the middle spheres. The fiber of $c$ is
\[
 C\simeq\operatorname{holim}(\cdots\xrightarrow pS\xrightarrow pS)
   \simeq F(S[1/p],S).
\]
Multiplication by $p$ is an equivalence on this function spectrum,
because it is an equivalence on its first argument. It remains an
equivalence after ordinary smash with $K(a)$. But
$\pi_*(K(a)\wedge C)$ is a module over $K(a)_*$ and is annihilated
by $p$. Thus it is zero, proving the first claim.

A map to a $K(t)$-local target consequently extends across $c$.
The target is also $L_r$-local when $t\leq r$, since an $L_r$-acyclic
spectrum is $K(t)$-acyclic. The extension therefore factors through
$L_rX$.
\end{proof}

\begin{proposition}[An actual sphere kernel]\label{prop:kernel-criterion}
Let $n\geq2$, let $1\leq t\leq n-1$, and let $z\in\pi_qS$.
Suppose a map $S\to Y$ to a $K(t)$-local spectrum carries $z$ to a
nonzero element, and the unit $S\to L_{K(n)}S$ carries $z$ to zero.
Then the image $\bar z\in\pi_qL_{n-1}X$ is nonzero and
$\pi_q(i_{n,X})(\bar z)=0$.
\end{proposition}
\begin{proof}
Lemma~\ref{lem:completion} factors the detecting map through
$L_{n-1}X$, so $\bar z$ is nonzero. Write $\eta$ for $K(n)$-localization
and $\lambda$ for $L_{n-1}$-localization. The relevant two squares are
\[
\begin{CD}
 S @>{c}>> X\\
 @V{\eta_S}VV @VV{\eta_X}V\\
 L_{K(n)}S @>{L_{K(n)}c}>> L_{K(n)}X
\end{CD}
\qquad
\begin{CD}
 X @>{\eta_X}>> L_{K(n)}X\\
 @V{\lambda_X}VV @VV{\lambda_{L_{K(n)}X}}V\\
 L_{n-1}X @>{i_{n,X}}>> L_{n-1}L_{K(n)}X.
\end{CD}
\]
Both commute by naturality of localization units. Equivalently,
\[
 \eta_Xc\simeq L_{K(n)}(c)\eta_S,
 \qquad i_{n,X}\lambda_X\simeq\lambda_{L_{K(n)}X}\eta_X.
\]
The first sends $z$ to zero, and the second then gives the asserted
kernel of the exact map \eqref{eq:weak-unit}. A retraction would induce
a left inverse on this homotopy group and is therefore impossible.
\end{proof}

\section{The named sphere product at the lower height}
\label{sec:lower-sphere-product}

The classical first beta element supplies the lower detector at every
odd prime. The fixed-point calculation identifies its image under an
actual unit, and the ordinary multiplicative filtration then detects
its long power. This gives one named sphere product for both upper
heights.

Fix an odd prime $p$ and put
\[
 h=p-1,\qquad d=2p(p-1)-2,\qquad
 z_p=\beta_1^{(p-1)^2},\qquad \beta_1\in\pi_dS.
\]

\begin{theorem}[The lower sphere product]\label{thm:lower-sphere-product}
The image of $z_p$ is nonzero in each of
\[
 \pi_{d(p-1)^2}L_{K(p-1)}S,\qquad
 \pi_{d(p-1)^2}L_{K(p-1)}X,\qquad
 \pi_{d(p-1)^2}L_rX\quad(r\geq p-1).
\]
All three images are induced by the sphere unit and, for the completed
input, by the completion map $S\to X$.
\end{theorem}

\subsection{The actual unit in the ordinary detector}

Let $E_h$ be Morava $E$-theory over $\F_{p^h}$. Write $\St_h$ for
the pure stabilizer and
$\G_h=\St_h\rtimes\Gal(\F_{p^h}/\F_p)$ for the extended stabilizer.
Choose the finite subgroup $G\subset\G_h$ whose intersection with the pure
stabilizer is $C_p\rtimes C_{h^2}$ and whose Galois projection is onto,
as in \cite[Section~2]{heard2015}. Put
\[
 Y=E_h^{hG},\qquad \eta:S\longrightarrow Y.
\]
Here $Y$ is the ordinary homotopy fixed point spectrum and $\eta$ is
its unit as a commutative ring spectrum. Devinatz--Hopkins identify the
ordinary finite-group fixed-point spectral sequence with the local
Adams spectral sequence and give strong convergence
\cite[Theorems~2(ii) and~3]{dh2004}. Its decreasing homotopy filtration
$F^s\pi_*Y$ is therefore separated, with
\[
 F^s\pi_qY/F^{s+1}\pi_qY\cong E_\infty^{s,q+s}.
\]

The Hopkins--Miller calculation gives the positive-filtration part of
its $E_2$ page as the corresponding part of
\[
 \F_p[\alpha,\beta,\Delta^{\pm1}]/(\alpha^2),\qquad
 |\alpha|=(1,2h),\quad |\beta|=(2,2ph),\quad
 |\Delta|=(0,2ph^2);
\]
see \cite[Proposition~5.5]{heard2015} and
\cite[Proposition~2.7]{hms2017}. We need two further assertions of that
calculation. First, Heard identifies $\beta$ with the image of the
named stable $\beta_1$ under $\eta$
\cite[proof of Theorem~5.6]{heard2015}. Allowing for the choice of
cohomology generator, this says
\begin{equation}\label{eq:lower-unit-detection}
 \eta_*(\beta_1)\in F^2\pi_dY,\qquad
 [\eta_*(\beta_1)]=c\beta\in E_\infty^{2,2ph},
 \qquad c\in\F_p^\times.
\end{equation}
The brackets denote the class modulo $F^3$. The Ext-to-group-cohomology
comparison and permanence are also recorded in the discussion after
\cite[Proposition~2.7]{hms2017}. Second, in positive even filtration
Heard's ordinary $E_\infty$ calculation is
\begin{equation}\label{eq:lower-even-infinity}
 E_\infty^{\mathrm{even},*}
   \cong \F_p[\Delta^{\pm p}][\beta]/(\beta^{h^2+1}),
 \qquad \text{in positive filtration},
\end{equation}
by \cite[Lemma~5.8]{heard2015}. In particular $\beta^{h^2}$ is
nonzero. These inputs hold for every odd prime; the named detection
is not restricted to $p=5$. They concern the extended group $G$ and
$\F_p$ coefficients. The pure subgroup and its larger coefficient
field will be used separately in the exponent argument.

\begin{lemma}[The ordinary filtered product]\label{lem:lower-bar-product}
Let a finite group $Q$ act coherently on a commutative ring spectrum $R$.
The ordinary homotopy fixed point spectral sequence has
\[
 E_2^{s,t}=H^s(Q;\pi_tR)
\]
and has the cup product on $E_2$ and the derivation rule for its
differentials. Whenever this spectral sequence converges to
$\pi_*R^{hQ}$ with a separated homotopy filtration whose successive
quotients are $E_\infty$, its product on $E_\infty$ is the
associated-graded product for that actual homotopy filtration.
\end{lemma}

The proof is given in Appendix~\ref{sec:ordinary-products}.

\begin{proposition}[The actual product in the detector]
\label{prop:lower-power}
The unit $\eta:S\to Y$ carries $z_p=\beta_1^{h^2}$ to a nonzero
class in $\pi_{dh^2}Y$.
\end{proposition}
\begin{proof}
Apply Lemma~\ref{lem:lower-bar-product} to the ordinary spectral
sequence just described. The class of $\eta_*(\beta_1)^{h^2}$ in
$F^{2h^2}\pi_{dh^2}Y/F^{2h^2+1}\pi_{dh^2}Y$ is
\[
 c^{h^2}\beta^{h^2}\ne0
\]
by \eqref{eq:lower-unit-detection} and
\eqref{eq:lower-even-infinity}. Hence the homotopy class itself is
nonzero. Since $\eta$ is a ring map, this power is exactly
$\eta_*(\beta_1^{h^2})$, the image of the specified sphere product.
\end{proof}

\begin{proof}[Proof of Theorem~\ref{thm:lower-sphere-product}]
The spectrum $Y$ is $K(h)$-local, since $E_h$ is local and local
objects are closed under homotopy limits. Its unit therefore factors
through $L_{K(h)}S$, where the image of $z_p$ must be nonzero by
Proposition~\ref{prop:lower-power}. Lemma~\ref{lem:completion} makes
$L_{K(h)}(S\to X)$ an equivalence; naturality carries that nonzero
image to $L_{K(h)}X$. For every $r\geq h$, the same lemma factors
the detecting unit through $S\to X\to L_rX\to Y$. Its nonzero value
on $z_p$ proves the remaining assertion.
\end{proof}

At $p=5$ this is the named class $\beta_1^{16}$ in degree
$38\cdot16=608$. Section~\ref{sec:kernels} kills its image at height
five as part of the uniform upper argument.

\section{Semilinear descent and an adic bound}\label{sec:adic}
At height $a=p$ for $p\geq5$, and at height $a=p+1$ for $p\geq7$,
the upper arguments need continuous cohomology to vanish above degree
$a^2+1$, for both sphere coefficients
and finite-test coefficients. We first bound the cohomology of finite
adic quotients and then pass to their countable inverse limit. At height
$p$ the Galois quotient has order $p$, and its semilinear action on Witt
coefficients supplies the required descent. At height $p+1$ the full
extended group has finite $p$-cohomological dimension.

Put $k_a=\F_{p^a}$ and write
\[
 (E_a)_*=O_a[u^{\pm1}],\quad
 O_a=W(k_a)[[u_1,\ldots,u_{a-1}]],\quad
 \m_a=(p,u_1,\ldots,u_{a-1}),\quad |u|=-2,
\]
and $\G_a=\St_a\rtimes\Gal(k_a/\F_p)$. Here $\St_a$ is the
pure stabilizer over $k_a$.

Exactness of Galois invariants for $p$-complete twisted Witt modules is
proved in \cite[Lemma~1.14]{goerss-hopkins2020}. The explicit contraction
below applies to arbitrary semilinear modules.

\begin{lemma}[Trace-one contraction]\label{lem:semilinear}
Let a finite group $Q$ act on a commutative ring $W$. Suppose
$w\in W$ satisfies $\sum_{g\in Q}g(w)=1$. For any semilinear
$W$-module $P$, one has $H^i(Q,P)=0$ for $i>0$, with no finiteness
or flatness hypothesis on $P$.
In particular this holds for $W=W(\F_{p^a})$ and
$Q=\Gal(\F_{p^a}/\F_p)$, without a prime-to-$p$ hypothesis on $a$.
\end{lemma}
\begin{proof}
Use homogeneous group cochains: a degree-$q$ cochain is an equivariant
function $f:Q^{q+1}\to P$, with differential the alternating deletion
of coordinates. For $q\geq1$ define
\[
 (hf)(g_0,\ldots,g_{q-1})
   =\sum_{\gamma\in Q}\gamma(w)f(\gamma,g_0,\ldots,g_{q-1}).
\]
For $g\in Q$, replacing $\gamma$ by $g\gamma$ proves that $hf$ is equivariant:
semilinearity converts $(g\gamma)(w)\,g f$ into
$g(\gamma(w)f)$. In $hd+dh$, the terms deleting one of the displayed
$g_i$ cancel in pairs. The term deleting the inserted coordinate is
$\sum_\gamma\gamma(w)f=f$. Thus $hd+dh=\mathrm{id}$ in positive
degrees, proving acyclicity for arbitrary $P$.

For the Witt-ring case, the finite-field trace is onto. Lift a residue
trace-one element to $w_0\in W$. Its Witt trace lies in $\Z_p$ and is
congruent to one modulo $p$, hence is a unit. Dividing $w_0$ by this
invariant unit supplies $w$. We have not divided by $|Q|$.
\end{proof}

\begin{lemma}[Finite quotients and the extra adic degree]\label{lem:adic}
Let $G=S_0\rtimes Q$ be profinite with $Q$ finite, and let $D\geq0$
be an integer such that $\cd_p(S_0)\leq D$ for discrete $p$-primary modules.
Let $\{M_j\}_j$ be a countable tower of finite discrete $p$-primary
$G$-modules with surjective transitions, and give
$M=\varprojlim_j M_j$ the inverse-limit topology. Suppose each $M_j$ is a
$W$-module, $S_0$ acts $W$-linearly, $Q$ acts semilinearly, and $W$ admits the
trace-one element of Lemma~\ref{lem:semilinear}. Then
\[
 H^s_{\cts}(G,M_j)=0\quad(s>D),\qquad
 H^s_{\cts}(G,M)=0\quad(s>D+1).
\]
The second conclusion also holds if the first displayed finite-quotient
vanishing is supplied directly, without a semidirect-product hypothesis.
\end{lemma}
\begin{proof}
The pure-subgroup cochain complex is $W$-linear and has semilinear
$Q$-action, including the conjugation action on its arguments: for an
inhomogeneous $b$-cochain,
\[
 (\gamma f)(s_1,\ldots,s_b)=
 \gamma\bigl(f(\gamma^{-1}s_1\gamma,\ldots,
                  \gamma^{-1}s_b\gamma)\bigr).
\]
The pure differential is $W$-linear and commutes with this action. Its
cohomology is therefore semilinear, whether or not that cohomology has
separately been shown finite. The continuous Hochschild--Serre sequence
at the discrete coefficient level is
\[
 H^i\bigl(Q,H^b_{\cts}(S_0,M_j)\bigr)
 \Longrightarrow H^{i+b}_{\cts}(G,M_j).
\]
Lemma~\ref{lem:semilinear} kills its positive $Q$-columns, giving
\[
 H^s_{\cts}(G,M_j)=H^s_{\cts}(S_0,M_j)^Q.
\]
This proves the first bound. Continuous cochains into $M$ are compatible
cochains into its finite quotients. Their transition maps are degreewise
surjective: a continuous cochain has a finite value set, and any choice
of lifts of those values gives a continuous cochain into the next finite
quotient. No equivariant section is needed for inhomogeneous cochains.
The countable inverse-limit sequence is therefore
\begin{equation}\label{eq:adic-milnor}
 0\longrightarrow\varprojlim_j{}^1 H^{s-1}_{\cts}(G,M_j)
 \longrightarrow H^s_{\cts}(G,M)
 \longrightarrow\varprojlim_j H^s_{\cts}(G,M_j)\longrightarrow0.
\end{equation}
Both outer terms vanish for $s>D+1$. At the remaining boundary degree,
the same sequence gives
\[
 H^{D+1}_{\cts}(G,M)
 \cong\varprojlim_j{}^1 H^D_{\cts}(G,M_j).
\]
Thus the bound retains the possible extra inverse-limit degree.
\end{proof}

\begin{proposition}[The two upper heights]\label{prop:cohomology}
Let $V_0$ be a finite ordinary $p$-local spectrum, let $DV_0=F(V_0,S)$,
and give $M=\pi_q(E_a\wedge DV_0)$ its $\m_a$-adic topology and
Morava action. Then
\[
 H^s_{\cts}(\G_a,M)=0\qquad(s>a^2+1)
\]
in each of the following cases: $a=p$ with $p\geq5$, and $a=p+1$
with $p\geq7$.
\end{proposition}
\begin{proof}
The pure stabilizer is the maximal-order unit group in the central
division algebra over $\Q_p$ of dimension $a^2$, so it is compact
$p$-adic analytic of that dimension \cite[Sections~2.1--2.2]{morava1985}.
It has no element $g$ of order $p$ in either stated range. Indeed,
$(g-1)\Phi_p(g)=0$ in a division algebra and $g\ne1$ imply
$\Phi_p(g)=0$. The polynomial $\Phi_p(X+1)$ is Eisenstein, so $g$
generates the degree-$(p-1)$ subfield $\Q_p(\zeta_p)$. The division
algebra would then be a vector space over this subfield, forcing
$p-1\mid a^2$. For $a=p$ this forces $p-1\mid1$; for $a=p+1$ it
forces $p-1\mid4$. These are impossible in the stated ranges.
The Lazard--Serre theorem therefore gives
$\cd_p(\St_a)=a^2$ for discrete $p$-primary coefficients
\cite[Section~1, Corollary~(1)]{serre1965}.

For each $q$, the coefficient module $M$ is finitely generated over the
complete Noetherian ring $O_a$, by the finite construction of $DV_0$.
It is complete and separated; its quotients $M_j=M/\m_a^jM$ are finite
discrete $p$-primary modules with surjective transitions. This uses no
flatness of $M$ and allows torsion and both parities. The maximal ideal
is invariant. For $V_0=S$, continuity of the action on these quotients
is part of the Morava action used by Devinatz--Hopkins
\cite[Theorem~1(ii) and Remark~1.3]{dh2004}. For every finite $V_0$,
the actual action and its continuity are also obtained by the ordinary
cooperation construction of Proposition~\ref{wup:cochains} below.

The pure stabilizer fixes $W(k_a)$. It fixes the residue field and thus
the unique Hensel lifts of the roots of $X^{p^a}-X$, which are the
Teichm\"uller representatives; continuity then fixes their convergent
Witt expansions. Hence its action on $M_j$ is Witt-linear, while
$Q=\Gal(k_a/\F_p)$ acts semilinearly by the Witt lift. At $a=p$,
Lemma~\ref{lem:adic} with $D=p^2$ gives the bound. The full group has
the split Galois subgroup of order $p$; this conclusion is specific to
the stated coefficients. The cohomological-dimension theorem was
applied only to $\St_p$.

At $a=p+1$ with $p\geq7$, the Galois quotient has order $p+1$, prime to $p$.
An order-$p$ element of $\G_a$ would therefore lie in $\St_a$, where
we have just excluded it. The full group is a finite extension of the
compact $p$-adic analytic pure group and has the same dimension $a^2$.
The same dimension theorem now applies to the full group itself and
gives $H^s_{\cts}(\G_a,M_j)=0$ for $s>a^2$. Apply only the
inverse-limit part of Lemma~\ref{lem:adic} to this finite-quotient
bound. In both cases the finite quotients vanish for $s>a^2$, and
\eqref{eq:adic-milnor} gives the asserted bound $s>a^2+1$ for $M$.
\end{proof}

For the sphere-kernel argument, take $V_0=S$, so that $M=(E_a)_q$.
The exponent argument uses the same bound for every finite $V_0$,
including modules with torsion and both parities. In both applications
the topology is the full maximal-ideal topology, as in the
continuous-cochain convention of Devinatz--Hopkins
\cite[Remark~1.3]{dh2004}.

\begin{remark}[A sharper line for sphere coefficients]\label{rem:sphere-sharp}
For the sphere homotopy fixed point spectral sequence, Bobkova and
collaborators state an $E_2$ vanishing line at $s=n^2+1$ when
$p-1\nmid n$ \cite[Equation~(2) and Remark~1.2]{bobkova2025}.
At $p=n=5$ this gives
$H^s_{\cts}(\G_5;(E_5)_q)=0$ for $s>25$.
That statement concerns the ordinary sphere coefficients $(E_n)_q$.
It does not assert the same line for every finite-test module
$\pi_q(E_n\wedge DV_0)$ or identify an actual image filtration.
Proposition~\ref{prop:cohomology} supplies the common conservative bound
used in this paper.
\end{remark}

\section{Upper annihilation and the canonical kernels}\label{sec:kernels}

At each upper height, torsion and parity place the unit image of
$\beta_1$ in Adams filtration two. The cohomological bound of
Section~\ref{sec:adic} makes the sufficiently high filtration zero.
To apply this to the same power detected at the lower height, we identify the Adams
filtration with images of tensor powers of the unit fiber. This also
proves its multiplicativity directly.

\subsection{Relative Adams stages}
\label{wimg:local-tower}

We first isolate the finite-stage comparison. In a stable symmetric
monoidal $\infty$-category with exact tensor, let $E$ be a unital algebra.
Call an object $E$-injective if it is a retract of $E\otimes Z$ for
some $Z$. A relative $E$-Adams resolution of $V$ consists of triangles
\[
 T_{s+1}\longrightarrow T_s\xrightarrow{q_s}Q_s,
 \qquad T_0=V,
\]
where each $Q_s$ is $E$-injective and $E\otimes q_s$ is split monic.
All stages map to $V$ by their composites down the tower. This is the
relative convention of \cite[Appendix~A]{dh2004}.

\begin{lemma}[Equality of relative Adams-stage images]
\label{lem:adams-stage-images}
Two such resolutions of $V$ have the same stage-$s$ image in $[W,V]$
for every object $W$ and every $s\geq0$.
\end{lemma}
\begin{proof}
Suppose a map $f_s:T_s\to T'_s$ over $V$ has been constructed, and
put $g_s=q'_sf_s$. Choose a retract
\[
 Q'_s\xrightarrow{\iota_s}E\otimes Z_s
       \xrightarrow{\rho_s}Q'_s,\qquad \rho_s\iota_s=1,
\]
and a splitting
$\sigma_s:E\otimes Q_s\to E\otimes T_s$ with
$\sigma_s(1_E\otimes q_s)=1$. The composite
\begin{equation}\label{wimg:extension}
\begin{aligned}
 \overline g_s:\quad Q_s
 &\xrightarrow{\eta_{Q_s}}E\otimes Q_s
 \xrightarrow{\sigma_s}E\otimes T_s\\
 &\xrightarrow{1_E\otimes(\iota_sg_s)}E\otimes E\otimes Z_s
 \xrightarrow{\mu\otimes1}E\otimes Z_s
 \xrightarrow{\rho_s}Q'_s
\end{aligned}
\end{equation}
extends $g_s$ along $q_s$. Indeed, unit naturality gives
$\eta_{Q_s}q_s=(1_E\otimes q_s)\eta_{T_s}$. The splitting removes
$1_E\otimes q_s$, and a second use of unit naturality followed by
the module unit law gives
\[
 \overline g_s q_s
 =\rho_s(\mu\otimes1)(1_E\otimes(\iota_sg_s))\eta_{T_s}
 =\rho_s\iota_sg_s=g_s.
\]
Complete this square to a map of fiber triangles to obtain
$f_{s+1}:T_{s+1}\to T'_{s+1}$. Starting with
$f_0=\mathrm{id}_V$ gives maps over $V$ at every finite stage,
and hence one inclusion between the stage-$s$ images in $[W,V]$.
Repeating the construction in the other direction proves equality.
The comparison begins at stage zero on $V$, so the indices agree.
Only these finite-stage maps are used; no equivalence of towers is
asserted.
\end{proof}

\subsection{The separated upper filtration}

For a positive height $a$, work in the local category
$\mathcal C_a=\mathrm{Sp}_{K(a)}$, with
\[
 B_a=L_{K(a)}S,\qquad
 U\otimes_a V=L_{K(a)}(U\wedge V),\qquad
 I_a=\fib(B_a\to E_a),\qquad \epsilon_a:I_a\to B_a.
\]
The tensor unit is $B_a$. These local tensors will be used to
describe homotopy classes of the underlying ordinary spectra.

\begin{lemma}[The actual Adams filtration]\label{wimg:descent}
The strongly convergent descent spectral sequence
\begin{equation}\label{wimg:morava-page}
 E_2^{s,q}=H^s_{\cts}(\G_a;(E_a)_q)
 \Longrightarrow\pi_{q-s}B_a
\end{equation}
has the separated homotopy filtration
\begin{equation}\label{wimg:actual-filtration}
 F^s\pi_*B_a
 =\im\bigl(\pi_*I_a^{\otimes_a s}
       \xrightarrow{(\epsilon_a^{\otimes_a s})_*}\pi_*B_a\bigr),
 \qquad I_a^{\otimes_a0}=B_a.
\end{equation}
Its successive quotients are $E_\infty^{s,*}$ with the usual total
degree, and $F^rF^s\subseteq F^{r+s}$.
\end{lemma}
\begin{proof}
Devinatz--Hopkins give \eqref{wimg:morava-page} with strong convergence
and identify it with the local $E_a$-Adams spectral sequence
\cite[Theorem~1(iii)--(iv)]{dh2004}; the full extended group is open
in itself. The coefficient topology is the full maximal-ideal-adic
topology of \cite[Remark~1.3]{dh2004}. Their ideal
$(p,v_1,\ldots,v_{a-1})$ agrees degreewise with
$\m_a=(p,u_1,\ldots,u_{a-1})$, since
$v_i=u_i u^{1-p^i}$ and $u$ is invertible.

The tensor powers $T_s=I_a^{\otimes_a s}$ form resolution triangles
\begin{equation}\label{wimg:power-triangles}
 T_{s+1}\longrightarrow T_s
     \xrightarrow{q_s}E_a\otimes_a T_s.
\end{equation}
Their right terms are $E_a$-injective, and multiplication on $E_a$
retracts $E_a\otimes_a q_s$. Thus they form a relative Adams
resolution in the convention just used. By
Lemma~\ref{lem:adams-stage-images}, their stage-$s$ images agree
with those in the Devinatz--Hopkins resolution at the same index.
This proves \eqref{wimg:actual-filtration}. In particular $T_0=B_a$
and $T_1=I_a$ give the full group and the kernel of
$\pi_*B_a\to\pi_*E_a$, respectively. Strong convergence supplies
separation and the stated associated-graded quotients for these
actual image groups.

A map from a suspended sphere to a local object is equivalently a
map from the suspended local unit. Tensoring two lifts through
$I_a^{\otimes_a r}$ and $I_a^{\otimes_a s}$ therefore gives a lift
of their product through $I_a^{\otimes_a(r+s)}$. This proves
multiplicativity. Monoidal localization carries the sphere unit
and its products to the local unit, so these are also the products
of the actual sphere images. Smashing localization is not needed.
\end{proof}

\begin{proposition}[The zero upper tail]\label{wimg:upper-filtration}
Let $a=p$ with $p\geq5$, or $a=p+1$ with $p\geq7$. Then
\begin{equation}\label{wimg:zero-tail}
 F^s\pi_*B_a=0\qquad(s\geq a^2+2).
\end{equation}
\end{proposition}
\begin{proof}
Proposition~\ref{prop:cohomology}, with the finite test $V_0=S$,
gives
\begin{equation}\label{wimg:coefficient-line}
 H^s_{\cts}(\G_a;(E_a)_q)=0
 \qquad(s>a^2+1,\ q\in\Z).
\end{equation}
At height $p$ this uses the pure stabilizer dimension and the
semilinear Witt trace-one contraction for the order-$p$ Galois
quotient. At height $p+1$ the full extended group has no $p$-torsion,
so its dimension gives the finite-quotient bound. Both cases retain
the possible adic inverse-limit contribution in degree $a^2+1$.

By Lemma~\ref{wimg:descent}, the successive quotients of the actual
filtration from $a^2+2$ onward are zero. That tail is consequently
constant, and separatedness makes its common value zero. This proves
\eqref{wimg:zero-tail}. It is a statement about homotopy images;
the exponent argument will require a separate proof of map nullity.
\end{proof}

\subsection{Annihilating the detected sphere product}
\label{wimg:kernel-assembly}

Put $d=2p(p-1)-2$ and $k=(p-1)^2$.

\begin{proposition}[Upper vanishing of the same power]
\label{wimg:upper-power}
For every $x\in\pi_dS$, the unit image of $x^k$ is zero in
$\pi_{dk}B_p$ when $p\geq5$. When $p\geq7$, the image of the
same product is also zero in $\pi_{dk}B_{p+1}$.
\end{proposition}
\begin{proof}
Let $a$ be either upper height in its stated range, and let
$x_a\in\pi_dB_a$ be the unit image of $x$. Positive stable sphere
stems are torsion: suspend into an odd-dimensional sphere in the stable
range and apply Serre's finiteness theorem
\cite[Chapter~V, Section~3, Proposition~3]{serre1951}.
Since $(E_a)_d$ is torsion-free, $x_a$ maps to zero in
$\pi_dE_a$ and belongs to $F^1\pi_dB_a$.

Since $d$ is even and $E_a$ has only even homotopy,
\[
 F^1\pi_dB_a/F^2\pi_dB_a
     \cong E_\infty^{1,d+1}=0.
\]
Thus $x_a\in F^2$, and multiplicativity in
Lemma~\ref{wimg:descent} gives $x_a^k\in F^{2k}\pi_{dk}B_a$.
The numerical margins above the last possible filtration are
\[
\begin{aligned}
 2(p-1)^2-(p^2+1)&=p^2-4p+1>0 &&(p\geq5),\\
 2(p-1)^2-\bigl((p+1)^2+1\bigr)&=p(p-6)>0 &&(p\geq7).
\end{aligned}
\]
Therefore $2k\geq a^2+2$, and \eqref{wimg:zero-tail} gives
$x_a^k=0$. This is the unit image of the original sphere product
$x^k$, by the product identification in Lemma~\ref{wimg:descent}.
\end{proof}

\begin{proof}[Proof of Theorems~\ref{thm:height-plus-one-kernel}
and~\ref{thm:height-p-kernel}, and Corollary~\ref{thm:five}]
Use the single named product $z_p=\beta_1^{(p-1)^2}$.
Proposition~\ref{prop:lower-power} detects it under the actual unit
to the $K(p-1)$-local spectrum $Y=E_{p-1}^{hG}$.
Proposition~\ref{wimg:upper-power}, applied to $x=\beta_1$, kills
that same product in $B_p$ for $p\geq5$ and in $B_{p+1}$ for
$p\geq7$.

Apply Proposition~\ref{prop:kernel-criterion} with $t=p-1$ and
$n=p$ or $n=p+1$, respectively. The result is a nonzero image of
$z_p$ in $\pi_{dk}L_{n-1}X$ killed by the actual canonical map
$i_{n,X}$. Its two naturality squares identify that map after
completion. A retraction would induce a left inverse on this
homotopy group, which is impossible.

For $p=5$ the product is $\beta_1^{16}$ in degree $608$, and its
upper image lies in $F^{32}\pi_{608}B_5$, while the zero tail
starts at $F^{27}$. Taking $n=5$ gives the named prime-five
statement as well.
\end{proof}

\section{Finite descent and exponent bounds}\label{sec:towers}
The exponent argument uses two consequences of a finite descent bound.
First, such a bound forces a horizontal vanishing line on a specified
page of the descent spectral sequence. Second, an $E_2$ vanishing line,
together with some finite descent bound, forces a specified fiber map
to have zero image after every chosen exact test. We prove these
statements directly from finite tensor cubes, keeping tested images
distinct from null maps.

\subsection{Finite tensor cubes}\label{par:tower-setup}
Let $\mathcal C$ be a stable symmetric monoidal $\infty$-category whose tensor is
exact in each variable. Let $A$ be an associative unital algebra, with unit
$\one$, and put $J=\fib(\one\to A)$ and $\epsilon:J\to\one$.
For $V\in\mathcal C$ form the augmented Amitsur object
\[
 V\longrightarrow C_A^\bullet(V),\qquad
 C_A^q(V)=A^{\otimes(q+1)}\otimes V.
\]
Cofaces insert the unit and codegeneracies multiply adjacent $A$-factors.
The partial totalization $\Tot_s$ is the limit over the full simplex
category on $[0],\ldots,[s]$, including codegeneracies. The cubical
comparison is standard
descent machinery \cite[Proposition~2.14, pp.~1007--1008]{mnn2017}; its proof fixes the
finite-stage indexing and functoriality used here.

\begin{lemma}[Full finite tensor cube]\label{lem:cube}
For $s\geq0$ there is a natural equivalence of arrows over $\mathrm{id}_V$
in $\mathcal C_{/V}$:
\[
 \begin{aligned}
 &\left(\fib\bigl(V\to\Tot_s C_A^\bullet(V)\bigr)\longrightarrow V\right)\\
 &\qquad\simeq\left(J^{\otimes(s+1)}\otimes V
 \xrightarrow{\epsilon^{\otimes(s+1)}\otimes\mathrm{id}_V}V\right).
 \end{aligned}
\]
The transition on the fibers from stage $s+1$ to stage $s$ applies
$\epsilon$ to the final $J$-factor. Any exact functor
$\Phi:\mathcal C\to\mathrm{Sp}$ preserves this finite fiber sequence.
\end{lemma}
\begin{proof}
Let $\mathcal P_{\ne\varnothing}([s])$ be the poset of nonempty
subsets of $\{0,\ldots,s\}$. The functor
\[
 \theta_s:\mathcal P_{\ne\varnothing}([s])\longrightarrow\Delta_{\leq s},
 \qquad S\longmapsto[|S|-1]
\]
uses the order-preserving injection induced by each inclusion of subsets.
At $[l]$, with $l\leq s$, the comma category $\theta_s/[l]$ consists of
nonempty subsets equipped with a weakly increasing map to $[l]$. It is
the poset of nonempty faces of the simplicial complex $P(s,l)$ whose
simplices are lists
\[
 (i_0,j_0),\ldots,(i_b,j_b),\quad
 0\leq i_0<\cdots<i_b\leq s,\quad
 0\leq j_0\leq\cdots\leq j_b\leq l.
\]
This complex is contractible. The base case $P(0,0)$ is a point.
Every simplex either omits the first-coordinate value $s$ or contains
exactly one vertex $(s,j)$. Begin inductively with the cone of vertex
$(s,l)$ on $P(s-1,l)$.
This cone is contractible even when $l=s$, when its base is not covered
by induction. For each $j<l$, attach the cone of vertex $(s,j)$ along
its base $P(s-1,j)$. This is its exact intersection with the previous
complex and is contractible by induction, since $j\leq s-1$.
The intersections are simplicial subcomplexes, so each attachment
preserves contractibility. The functor is therefore
homotopy initial.

The partial totalization is consequently the limit of the punctured
$(s+1)$-cube obtained by tensoring copies of $\one\to A$ and then
tensoring with $V$. Its omitted vertex is $V$. Taking successive fibers
and using exactness of tensor gives $J^{\otimes(s+1)}\otimes V$,
with its canonical arrow to $V$. Removing the last direction of the
cube gives the stated transition. This is a genuinely finite cubical
limit, so any exact $\Phi$ preserves it. Finiteness of the list of
objects in $\Delta_{\leq s}$ alone would not justify that assertion.
\end{proof}

Fix an exact functor $\Phi:\mathcal C\to\mathrm{Sp}$. Put
$V'=\Phi V$, $U_l=\Phi(J^{\otimes l}\otimes V)$, and
$P_s=\Tot_s(\Phi C_A^\bullet(V))$ for $s\geq0$; set $P_s=0$
for $s<0$. Thus $U_0=V'$ and
\begin{equation}\label{eq:finite-fiber}
 U_{s+1}\longrightarrow V'\longrightarrow P_s
\end{equation}
is a compatible family of fiber sequences. Define its augmentation
filtration by
\begin{equation}\label{eq:augmentation-filtration}
 F^s\pi_dV'=\ker(\pi_dV'\to\pi_dP_{s-1})
           =\im(\pi_dU_s\to\pi_dV'),\qquad s\geq0.
\end{equation}
These groups are defined from finite stages. The argument below will
use them without identifying $V'$ with an infinite totalization.
Our spectral sequence convention is
\[
 E_2^{s,q}=H^s\pi_q(\Phi C_A^\bullet(V)),\qquad
 d_r:(s,q)\longmapsto(s+r,q+r-1).
\]
The full tower has normalized cosimplicial cochains on $E_1$, giving the
displayed $E_2$ page.

\subsection{Exponents and exact finite-page bounds}\label{sec:exponent-pages}
For $b\geq1$ say that $V$ has $A$-exponent at most $b$ if the actual map
\[
 e_b(V)=\epsilon^{\otimes b}\otimes\mathrm{id}_V:
 J^{\otimes b}\otimes V\longrightarrow V
\]
is null. This is stronger than vanishing on homotopy groups after a test.

\begin{lemma}[Exponent calculus]\label{lem:exponent-calculus}
An $A$-module has exponent at most one. Bounds are preserved by retracts
and by tensoring with arbitrary objects. For a cofiber sequence
$U\to V\to W$, bounds $b$ and $c$ on $U$ and $W$ give bound $b+c$ on $V$.
If $e_m(V)$ is null, $V$ is a retract of an object with $m$ successive
layers $A\otimes J^{\otimes i}\otimes V$, $0\leq i<m$.
\end{lemma}
\begin{proof}
The unit map of a module splits by its action, making its fiber map
null. Naturality gives retract invariance, and tensoring a null map
gives tensor invariance. In the cofiber sequence, the composite
$J^{\otimes c}\otimes V\to V\to W$ is null, so choose a lift
$\ell:J^{\otimes c}\otimes V\to U$. After precomposition by $b$
more $\epsilon$ factors, naturality identifies the resulting map into $U$
with $e_b(U)(1\otimes\ell)$, which is null. Its composite to $V$ is
$e_{b+c}(V)$. Finally the consecutive maps
$J^{\otimes(i+1)}\otimes V\to J^{\otimes i}\otimes V$ have the
stated cofibers. Their composites to $V$ filter $\cofib(e_m(V))$ by
these $m$ layers. Nullity of $e_m(V)$ makes $V$ a retract of that cofiber.
\end{proof}

These bounds describe the thick tensor ideal generated by $A$;
compare \cite[Definition~3.18]{mathew2016}. We now fix
the page convention for the quantitative use of nilpotence
\cite{mnn2017,mnn2019}.

\begin{proposition}[Pages and tested images]\label{prop:pages}
For any exact $\Phi:\mathcal C\to\mathrm{Sp}$ and integer $b\geq1$:
\begin{enumerate}
\item If $V$ has $A$-exponent at most $b$, then $E_{b+1}^{s,q}=0$ for all
      $s\geq b$ and $q$.
\item If $V$ has some finite $A$-exponent and $E_2^{s,q}=0$ for all
      $s\geq b$ and $q$, then $\pi_d\Phi(e_b(V))=0$ for every $d$.
\end{enumerate}
The second assertion is zero tested image, not nullity of a map.
\end{proposition}
\begin{proof}
In coordinates $(s,d)$ with $d=q-s$, the derived couple has
\[
 D_r^{s,d}=\im(\pi_dP_{s+r-1}\to\pi_dP_s)
\]
and maps
\[
 i:D_r^{s,d}\to D_r^{s-1,d},\qquad
 j:D_r^{s,d}\to E_r^{s+r,d-1},\qquad
 k:E_r^{s,d}\to D_r^{s,d}.
\]
Here the second $E$-index denotes total degree; no convergence claim is used.
Let $c$ be an actual exponent bound. Every length-$c$ transition in the
fiber tower $U_l$ is null, being $\Phi$ applied to a tensor multiple
of $e_c(V)$. For $s\geq0$ the boundary square from
\eqref{eq:finite-fiber} is
\[
 \begin{array}{ccc}
 \pi_dP_{s+c}&\longrightarrow&\pi_{d-1}U_{s+c+1}\\
 \downarrow&&\downarrow\,0\\
 \pi_dP_s&\longrightarrow&\pi_{d-1}U_{s+1}.
 \end{array}
\]
The image from $P_{s+c}$ has zero boundary and lifts to $\pi_dV'$;
conversely an augmentation class comes from $P_{s+c}$. Hence
\begin{equation}\label{eq:stable-derived-image}
 D_{c+1}^{s,d}=\im(\pi_dV'\to\pi_dP_s).
\end{equation}
Both sides are zero for negative $s$. These image groups have surjective
transition $i$. The augmentation into $P_s$ is injective on homotopy for
$s\geq c-1$, since its kernel is the image of a map factoring through
$e_c(V)$. Thus $i$ is an isomorphism for $s\geq c$. Exactness of the
derived couple gives $j=0$ and, for $s\geq0$, identifies
\[
 E_{c+1}^{s,d}\cong\ker(D_{c+1}^{s,d}\to D_{c+1}^{s-1,d})
                   \cong F^s\pi_dV'/F^{s+1}\pi_dV'.
\]
In particular this is zero for $s\geq c$. Taking $c=b$ proves (i).
For (ii) retain any $c$. The $E_2$ line makes these same quotients zero for
$s\geq b$. Since $F^c=0$, descending from $c$, or monotonicity when
$c<b$, gives $F^b=0$. Equation~\eqref{eq:augmentation-filtration}
is exactly the asserted zero image. This argument uses only finite
cubes and never assumes that $\Phi$ preserves an infinite limit.
\end{proof}

\section{A finite-page lower exponent bound}\label{sec:exponent-proof}
\label{sec:lower-exponent}
Fix a prime $p\geq5$ and put $t=p-1$, $h=p$, and $m=p^2+2$.
The lower exponent inequality in Theorem~\ref{thm:exponents} will follow
from a class in filtration greater than $m$ on page $E_{m+1}$.
Proposition~\ref{prop:pages}(i) excludes such a class whenever the
exponent is at most $m$. This argument uses the pure stabilizer subgroup
and its residue field, which we specify independently of the extended
group used for the sphere kernel.

\begin{theorem}[Pure-stabilizer cohomology]\label{thm:lower-hm}
For every odd prime $p$, put $t=p-1$ and let $E_t$ be Morava $E$-theory
for the height-$t$ Honda formal group over $k_t=\F_{p^t}$, with its
coherent stabilizer action. The pure stabilizer $\St_t$ contains a finite subgroup
$H\cong C_p\rtimes C_{t^2}$ for which
\begin{equation}\label{eq:lower-algebra}
 \begin{gathered}
 \widehat H^*(H;(E_t)_*)
 =k_t[\alpha,\beta^{\pm1},\Delta^{\pm1}]/(\alpha^2),\\
 |\alpha|=(1,2t),\qquad |\beta|=(2,2pt),\qquad
 |\Delta|=(0,2pt^2).
 \end{gathered}
\end{equation}
The ordinary-to-Tate comparison is multiplicative and identifies
$H^s(H;(E_t)_q)$ with the displayed Tate group for every $s>0$ and $q$.
\end{theorem}

The coherent action specializes to the prescribed Honda stabilizer
action \cite[Theorem~1(ii) and Remark~2.4]{dh2004}.
This is the subgroup denoted $F$ in
\cite[Equation~(1.2) and Theorem~4.5]{bbs2020}, where the height is
called $n$. The positive-filtration comparison is
\cite[Remark~2.2]{bbs2020}; the pure-group calculation before taking
Galois invariants also appears in
\cite[proof of Proposition~5.5]{heard2015}.
We use the cohomology and its bidegrees here. No differential from the
Tate spectral sequence is needed for the argument below.

Work in $\mathcal D_H=\Fun(BH,\Sp)$ with pointwise ordinary smash,
and put $A_H=F(H_+,S)$ with its translation action. For the exact
functor $(-)^{hH}$, the $A_H$-Amitsur spectral sequence is the ordinary
homotopy fixed point spectral sequence: its terms are functions on the
free $H$-sets $H^{q+1}$, and taking fixed points gives homogeneous group
cochains. Lemma~\ref{lem:lower-bar-product} supplies its cup product and
derivation rule without requiring convergence. We will apply the
finite-page exponent criterion to this tower.

\begin{proposition}\label{prop:lower-exponent}
The $H$-exponent of $E_t$ is greater than $m$.
\end{proposition}
\begin{proof}
We have $m<2t^2+1$, since the difference is $p^2-4p+1>0$.
Choose a positive $N$ divisible by $p$ with $2N>m$. We show that
$\beta^N\ne0$ on $E_{m+1}$ in bidegree $(2N,2ptN)$.
For $2\leq r\leq m$ all possible source and target filtrations are
positive, so \eqref{eq:lower-algebra} applies. Even $r$ gives odd
internal degree and a zero group. For odd $r$, internal-degree divisibility
forces $r=2ta+1$ with $1\leq a<t$.

An outgoing target must be $\alpha\beta^{N+ta}\Delta^l$.
Equality of internal degrees gives $a-1=pt(a+l)$, which forces $a=1$.
An incoming source must be $\alpha\beta^{N-ta-1}\Delta^l$,
whose internal-degree equation is $a-t=pt(a-l)$, impossible for
$1\leq a<t$. Thus only the outgoing differential $d_{2t+1}$ remains.
The same degree argument shows that $\beta$ exists up to that page;
its only possible incoming differential has length two and odd internal
source degree. The derivation rule on that page gives
\[
 d_{2t+1}(\beta^N)=N\beta^{N-1}d_{2t+1}(\beta)=0,
\]
because $\beta$ has even total degree, $p$ divides $N$, and the target is
$p$-torsion. With all incoming differentials and all other outgoing
possibilities excluded through page $m$, this class is nonzero on
$E_{m+1}$. Its filtration is $2N>m$, contradicting
Proposition~\ref{prop:pages}(i) if the exponent were at most $m$.
\end{proof}

At $p=5$, where $m=27$, one may take $N=15$: the class has filtration
$30$ on $E_{28}$ and obstructs exponent $27$. The proof needs only
this finite-page survival; it does not use a sphere representative or
the abutment of the ordinary spectral sequence.

\section{Ordinary cooperations and finite tests}\label{sec:cooperations}

The upper exponent bound requires the cohomological line of
Section~\ref{sec:adic} after every finite ordinary test.
We therefore identify the tested Amitsur coefficient complex, including
its actual coface and codegeneracy maps.
Fix a positive height $a$, and write
\[
 \begin{gathered}
 E=E_a,\qquad k_a=\F_{p^a},\qquad
 O=O_a=W(k_a)[[u_1,\ldots,u_{a-1}]],\\
 O_*=(E_a)_*=O[u^{\pm1}],\quad |u|=-2,\qquad
 \m=\m_a=(p,u_1,\ldots,u_{a-1}).
 \end{gathered}
\]
Use the coherent action of $\G_a=\St_a\rtimes\Gal(k_a/\F_p)$ by
commutative ring automorphisms, whose special-fiber action is the
prescribed Honda stabilizer action
\cite[Section~7, Proposition~7.3, Equation~(7.3), and
Corollaries~7.6--7.7]{gh2004}.
The local category $\mathcal C_a=\mathrm{Sp}_{K(a)}$ has unit
$B_a=L_{K(a)}S$ and tensor
$U\otimes_a V=L_{K(a)}(U\wedge V)$.
All unadorned smash products and function spectra below are ordinary.

\begin{proposition}[The finite-test coefficient complex]\label{wup:cochains}
For \(r\geq0\), let
\[
 Z_r=L_{K(a)}(E^{\wedge(r+1)}).
\]
For every finite ordinary \(p\)-local spectrum \(W\), there are
natural graded isomorphisms
\begin{equation}\label{wup:tested-functions}
 \pi_*(Z_r\wedge W)
 \cong C_{\cts}(\G_a^r,\pi_*(E\wedge W)).
\end{equation}
The first \(E\)-factor acts by constant coefficients.
The two-factor evaluation is multiplication after \(1\wedge g\).
In the cumulative coordinates \(1,g_1,g_1g_2,\ldots,g_1\cdots g_r\)
on the factors, these isomorphisms identify every coface and
codegeneracy of the unit Amitsur object with the continuous
inhomogeneous group-cochain operator.

Consequently, apply the exact ordinary functor \(F(V_0,-)\), for a
finite ordinary \(p\)-local spectrum \(V_0\), to the full unit
Amitsur tower for \(E_a\) in \(\mathcal C_a\). In the page convention
of Proposition~\ref{prop:pages}, its second page is
\begin{equation}\label{wup:cochain-page}
 E_2^{s,q}
 =H^s_{\cts}\bigl(\G_a,\pi_q(E_a\wedge D V_0)\bigr),
 \qquad D V_0=F(V_0,S).
\end{equation}
Here \(H^s_{\cts}\) means the cohomology of continuous cochains with
the adic topology on the indicated coefficient module.
\end{proposition}

We prove the comparison by first forming ordinary cooperations, then
completing their flat coefficient modules, and finally applying the
finite test. The next three lemmas justify these steps. Mixed
cooperations and their completion will also be used in
Section~\ref{sec:coherent-trace}.
The exact-homology input originates in Landweber's work
\cite{landweber1976}; we use the precise cooperation formulations cited
below, with $BP$ the $p$-local Brown--Peterson spectrum.

\begin{lemma}[Ordinary mixed cooperations]\label{wup:mixed-flat}
For positive heights \(a,b\), the ordinary smash product has even
homotopy
\[
 \pi_*(E_a\wedge E_b)\cong
 (E_a)_*\otimes_{BP_*}BP_*BP\otimes_{BP_*}(E_b)_*.
\]
It is flat over \((E_a)_*\) through the first factor.
\end{lemma}

\begin{proof}
Apply the ordinary Landweber cooperation calculation to the $p$-typical
flat Hopf algebroid $(BP_*,BP_*BP)$. This gives the displayed $BP_*$
formula. It is the $p$-typical form of the $MU_*$ formulas displayed in
\cite[Section~15, Proposition~15.1 and Corollary~15.4]{rezk1998} and
\cite[Lemma~4.4]{hovey2004}; Morava $E$-theory over $BP$ is included in
\cite[Example~0.1(d)]{hs-comodules2003}.
The two tensor products use the two different unit maps into \(BP_*BP\).
For a Landweber-exact \(BP_*\)-algebra \(C\), the module
\(BP_*BP\otimes_{BP_*}C\) is flat over the other \(BP_*\)-action.
Indeed, tensoring with it factors as the exact cofree-comodule functor
followed by the exact Landweber functor on comodules. This is also the
opposite-unit criterion of
\cite[Lemma~2.2]{hs-comodules2003}, with conjugation exchanging the
units when necessary. Base change along \(BP_*\to(E_a)_*\) proves the
asserted flatness. The displayed coefficient and cooperation rings are
even graded, which proves evenness. In particular this argument does
not require \((E_b)_*\) to be flat over \(BP_*\).
\end{proof}

The ring \(O\) is complete, Noetherian, regular, and local, with finite
residue field. If \(W\) is a finite ordinary \(p\)-local spectrum, then
\[
 M_*=\pi_*(E\wedge W)
\]
is finitely generated over \(O_*\). This follows by finite cofiber
constructions and retracts from the assertion for the sphere.
Periodicity identifies graded module questions with the two parities
over \(O\). Regularity therefore supplies a finite resolution
\cite[Tag~00O7]{stacks} of
\(M_*\) by finite graded free \(O_*\)-modules. Neither parity of
\(M_*\) is required to vanish.

\begin{lemma}[Base change and continuous functions]\label{wup:base-change}
Let \(D_0\) be an ordinary right \(E\)-module for which
\(\pi_*D_0\otimes_{O_*}(-)\) is exact on finitely generated graded
\(O_*\)-modules. The natural comparison
\begin{equation}\label{wup:relative-base-change}
 \pi_*D_0\otimes_{O_*}\pi_*(E\wedge W)
 \longrightarrow \pi_*(D_0\wedge W)
\end{equation}
is an isomorphism for every ordinary \(p\)-local spectrum \(W\).

For every profinite set \(T\), the functor \(C_{\cts}(T,-)\) is exact
on finitely generated \(O\)-modules with their \(\m\)-adic topologies.
For every finitely generated graded \(O_*\)-module \(M_*\), the
natural map is an isomorphism
\begin{equation}\label{wup:function-tensor}
 C_{\cts}(T,O_*)\otimes_{O_*}M_*
 \xrightarrow{\ \cong\ } C_{\cts}(T,M_*).
\end{equation}
The functions and the displayed graded isomorphism are understood
degreewise.
\end{lemma}

\begin{proof}
First suppose \(W\) finite. The equality
\[
 D_0\wedge W\simeq D_0\wedge_E(E\wedge W)
\]
is an equality of ordinary derived constructions. The comparison is
the balanced homotopy cross product for this relative smash, so it is
natural in both variables before choosing any resolution.
Realize a finite graded free surjection onto \(\pi_*(E\wedge W)\) by an \(E\)-module
map from a finite wedge of suspensions of \(E\). Surjectivity makes the
homotopy of its fiber the first syzygy. Repeat along a finite free
resolution. At the last stage, lifts of a free basis give an
equivalence from a finite free \(E\)-module, since they give an
isomorphism on all homotopy groups. Relative smash with \(D_0\) is
exact. Induction back through these fiber sequences, using the
assumed tensor exactness on each finitely generated syzygy, proves
\eqref{wup:relative-base-change} for finite \(W\).

Every ordinary \(p\)-local spectrum is a filtered homotopy colimit of
finite ordinary \(p\)-local spectra. Both sides of
\eqref{wup:relative-base-change} preserve those colimits: ordinary
smash and homotopy groups do so, as does tensoring a module with
\(\pi_*D_0\). Naturality extends the comparison to arbitrary \(W\).

For the assertion about functions, let \(M\to N\) be a surjection of
finitely generated \(O\)-modules. For every \(j\geq1\) the map
\[
 M/\m^{j+1}M\longrightarrow
 (M/\m^jM)\mathbin{\times}_{N/\m^jN}(N/\m^{j+1}N)
\]
is surjective. To see this, lift the first component to \(M\); its
error against the second component lies in \(\m^jN\), and this error
lifts from \(\m^jM\). Each quotient is a finite set. Given a continuous
map \(T\to N\), first lift its reduction modulo \(\m\), and then use
these surjections to lift successively while preserving the previous
choice. A lift on each finite value set is continuous. The compatible
system defines a continuous lift to \(M\), since finitely generated
modules are complete. Kernels carry their subspace topologies:
the relevant modules are compact Hausdorff, and the induced
continuous bijection from a kernel to its image is a homeomorphism.
It follows that \(C_{\cts}(T,-)\) is exact on these modules. Apply this
exact functor to a finite presentation of \(M\), and compare with
tensoring the same presentation by \(C_{\cts}(T,O)\). The comparison
is an isomorphism on the finite free terms and hence on \(M\).
Applying this argument in each parity proves
\eqref{wup:function-tensor}. In particular
\(C_{\cts}(T,O_*)\otimes_{O_*}(-)\) has precisely the tensor exactness
used in the first assertion.
\end{proof}

\begin{lemma}[Flat completion and reduction]\label{wup:completion-reduction}
Let \(N\) be an ordinary \(E\)-module with flat graded homotopy
\(U_*=\pi_*N\). Its actual \(K(a)\)-localization map induces
\[
 U_*\longrightarrow\pi_*L_{K(a)}N
 \cong\widehat{U_*}:=\varprojlim_j U_*/\m^jU_*.
\]
For a flat \(O\)-module \(U\) and every \(j\geq1\),
\[
 \widehat U/\m^j\widehat U\cong U/\m^jU.
\]
Both assertions apply degreewise to the periodic graded modules above.
\end{lemma}

\begin{proof}
The first assertion is the localization theorem for ordinary Morava
\(E\)-modules: \(K(a)\)-localization is derived \(\m\)-completion,
and flat homotopy has no higher derived completion. The natural
localization map induces the ordinary completion map
\cite[Propositions~3.4, 3.6 and Corollary~3.8]{rezk2009}.

Here is the algebraic reduction argument, including its needed
inverse-limit step. Choose a resolution \(P_\bullet\to N_0\) of a
finitely generated \(O\)-module \(N_0\) by finite free modules.
The tower \(P_\bullet\otimes_O U/\m^lU\) is degreewise surjective.
By flatness of \(U\), its first homology is
\[
 \operatorname{Tor}_1^O(N_0,O/\m^l)\otimes_O U.
\]
This homology tower is pro-zero. Indeed, for a finite first
presentation \(0\to K\to P_0\to N_0\to0\), the corresponding Tor group
is \((K\cap\m^lP_0)/\m^lK\). Artin--Rees
\cite[Tag~00IN]{stacks} gives a \(d\) such that
\(K\cap\m^lP_0\subseteq\m^{l-d}K\) for \(l\geq d\).
For any fixed \(j\), a transition from \(l\geq j+d\) to \(j\)
is therefore zero. Tensoring these zero transitions with \(U\)
preserves this conclusion.

The homology exact sequence for the countable inverse limit of this
degreewise-surjective tower now gives
\[
 N_0\otimes_O\widehat U
 \cong\varprojlim_l(N_0\otimes_O U/\m^lU);
\]
the possible \(\varprojlim^1\) of the first homology is zero because
that tower is pro-zero. Finite freeness of the \(P_i\) identifies
their tensor with \(\widehat U\) with their inverse-limit tensor.
Take \(N_0=O/\m^j\); for \(l\geq j\) the modules on the right are
\(U/\m^jU\). This proves the reduction formula.
\end{proof}

\begin{proof}[Proof of Proposition~\ref{wup:cochains}]
The completed cooperation theorem gives
\begin{equation}\label{wup:completed-cooperation}
 \pi_*L_{K(a)}(E\wedge E)
 \cong C_{\cts}(\G_a,O_*).
\end{equation}
At \(g\), the map is \(\mu(1\wedge g)\); the first factor acts by
constant scalars. This is the evaluation map of
\cite[Section~1 and Theorems~4.11 and 6.8]{hovey2004}.
The equivalent convention in
\cite[Proposition~2.2, Equation~(2.3), and Equations~(2.5)--(2.7)]{dh2004}
uses \(\mu(g^{-1}\wedge1)\); exchanging the factors and inverting
the group coordinate gives the convention above. The coefficient
topology is the adic one
\cite[Remark~1.3]{dh2004}.

Put \(P_*=\pi_*(E\wedge E)\), regarded as an \(O_*\)-module through
the first factor. Lemma~\ref{wup:mixed-flat} makes \(P_*\) flat.
The actual localization map, the reduction in
Lemma~\ref{wup:completion-reduction}, and
\eqref{wup:function-tensor} applied to \(O_*/\m^jO_*\) give
\begin{equation}\label{wup:ordinary-reductions}
 P_*/\m^jP_*
 \cong C_{\cts}(\G_a,O_*/\m^jO_*).
\end{equation}
This isomorphism is induced by the actual ordinary evaluation
\(E\wedge E\xrightarrow{1\wedge g}E\wedge E\xrightarrow{\mu}E\)
followed by reduction. It is not just an abstract isomorphism of
the two reduction modules.

To form higher powers, use a common first factor in the ordinary
relative smash products
\[
 (E_0\wedge E_1)\wedge_{E_0}\cdots
       \wedge_{E_0}(E_0\wedge E_r)
 \simeq E_0\wedge E_1\wedge\cdots\wedge E_r.
\]
Each \(E_i\) denotes the same spectrum \(E\); the subscripts mark
positions. Iterating \eqref{wup:relative-base-change}, including
its extension to arbitrary ordinary \(W\), identifies the homotopy
of this ordinary smash with \(P_*^{\otimes_{O_*}r}\).
The comparison inserts each pair into the first and one other
position and multiplies at the common position. These modules are
flat, since \(P_*\) is flat.

In each degree the coefficient ring \(O_*/\m^jO_*\) is finite.
Continuous functions to it are unions of function modules on finite
quotient sets of \(\G_a\). Tensors of the function modules on finite
sets are the function modules on their products, and every
continuous function on \(\G_a^r\) to a finite set factors through
finite quotients in the coordinates. It follows from
\eqref{wup:ordinary-reductions} that the reductions of
\(P_*^{\otimes_{O_*}r}\) are
\(C_{\cts}(\G_a^r,O_*/\m^jO_*)\). Apply
Lemma~\ref{wup:completion-reduction} only now, to this ordinary
\((r+1)\)-fold smash. The result is
\begin{equation}\label{wup:power-functions}
 \pi_*Z_r\cong C_{\cts}(\G_a^r,O_*).
\end{equation}
For \(r=0\) this is the locality of \(E\).

The maps in \eqref{wup:power-functions} remain the evaluations on
the non-first factors followed by multiplication. In fact each
ordinary evaluation extends through localization because its target
\(E\) is local. On homotopy its first-factor \(O_*\)-linearity makes
it adically continuous. Lemma~\ref{wup:completion-reduction}
makes the ordinary image dense, and the reductions above show that
the extension agrees with the displayed continuous function.
Thus the construction preserves the evaluation maps as well as the
coefficient modules.

For finite \(W\), the ordinary \(E\)-module \(Z_r\) satisfies the
exactness assumption in Lemma~\ref{wup:base-change}, by
\eqref{wup:power-functions} and \eqref{wup:function-tensor}.
Those comparisons give \eqref{wup:tested-functions}. This step
allows any finitely generated \(\pi_*(E\wedge W)\), including
torsion and both parities; no flatness of that module has been used.

The terms of the local unit Amitsur object are exactly \(Z_r\).
For finite \(V_0\), ordinary duality identifies
\(F(V_0,Z_r)\simeq Z_r\wedge D V_0\), so the previous formula
computes each cosimplicial homotopy group. Relabel the independent
evaluations by their cumulative coordinates
\[
 1,\quad g_1,\quad g_1g_2,\quad\ldots,\quad g_1\cdots g_r.
\]
The cochain identities below can be checked on ordinary evaluations before
localization. For finite $W$, their target $E\wedge W$ is $K(a)$-local,
since finite dualizable smash preserves local objects. They therefore
extend through localization and remain identities on the tested
coefficients. Insertion of the first unit acts on the coefficient; an
interior unit combines two consecutive coordinates, and
insertion of the last unit drops the last coordinate. Explicitly,
for a degree-\(r\) cochain \(c\),
\[
\begin{aligned}
 (\delta^0c)(g_1,\ldots,g_{r+1})
   &=g_1c(g_2,\ldots,g_{r+1}),\\
 (\delta^ic)(g_1,\ldots,g_{r+1})
   &=c(g_1,\ldots,g_i g_{i+1},\ldots,g_{r+1})
       &&(1\leq i\leq r),\\
 (\delta^{r+1}c)(g_1,\ldots,g_{r+1})
   &=c(g_1,\ldots,g_r).
\end{aligned}
\]
Multiplication of adjacent factors gives each codegeneracy, which
inserts the identity in the corresponding coordinate. These are all
the inhomogeneous cochain operators, including the degeneracies of
the full cosimplicial object. The first unit insertion also displays
the orbit function of every coefficient. Those functions are
continuous, and reduction gives continuous actions on the finite
invariant quotients by powers of \(\m_a\). The \(E_2\)-page of the
tower of finite partial totalizations is cosimplicial cohomology
of these homotopy groups, which is \eqref{wup:cochain-page}.
\end{proof}

\section{From finite tests to an actual null map}\label{sec:upper-nullity}

We now bound the exponent of the upper local unit by proving nullity
of a specified tensor-power map.

\begin{proposition}[The upper null map and its finite layers]
\label{wup:upper-null}
Let \(p\geq5\) be prime and put \(h=p\). In \(\mathcal C_h\), let
\(I_h=\fib(B_h\to E_h)\). The actual map
\[
 e_m(B_h):I_h^{\otimes_h m}\longrightarrow B_h,
 \qquad m=h^2+2=p^2+2,
\]
is null. Thus \(\expn_{E_h}(B_h)\leq m\), and \(B_h\) is a retract
of an object with \(m\) successive layers
\[
 E_h\otimes_h I_h^{\otimes_h i},\qquad 0\leq i<m.
\]
\end{proposition}

The coefficient line will make this map zero after every finite ordinary
precomposition, once some finite exponent is known. The following
property of its target then turns that tested vanishing into nullity.
Its relation to the Brown--Comenetz phantom criterion is
\cite[Proposition~4.11]{christensen-strickland1998}.

\begin{lemma}[Ordinary phantoms into a local sphere]\label{wup:phantom}
Let \(a\geq1\). An ordinary map \(f:U\to B_a\) of \(p\)-local
spectra is null if its precomposition with every map from a finite
ordinary \(p\)-local spectrum to \(U\) is null.
\end{lemma}

\begin{proof}
Let \(M_aS=\fib(L_aS\to L_{a-1}S)\), and let \(I_{\mathrm{BC}}\)
be the ordinary Brown--Comenetz dualizing spectrum for the injective
group \(\Q/\Z_{(p)}\). Gross--Hopkins duality, announced in
\cite[Theorem~6]{gross-hopkins1994} and proved in the form used here
in \cite[Theorem~10.2(b),(e)]{hs1999}, says that
\[
 D_a=F(M_aS,I_{\mathrm{BC}})
\]
is \(K(a)\)-local and invertible for \(\otimes_a\). The
identification with this monochromatic object uses
\cite[Theorem~6.19]{hs1999}.

For a local target \(Z\), the ordinary spectrum \(F(R,Z)\) is
local. A map into it from a \(K(a)\)-acyclic spectrum \(V\)
corresponds to a map \(V\wedge R\to Z\), and \(V\wedge R\) is
still acyclic. Hence ordinary function spectra compute the local
internal Hom, as in \cite[Theorem~7.1]{hs1999}. Invertibility of
\(D_a\), followed by ordinary closed adjunction, gives equivalences
of ordinary spectra
\[
 B_a\simeq F(D_a,D_a)
       \simeq F(D_a\wedge M_aS,I_{\mathrm{BC}}).
\]
The displayed smash is ordinary. Put \(C=D_a\wedge M_aS\).
Brown--Comenetz representability now gives, naturally in \(U\),
\[
 [U,B_a]\cong
 \operatorname{Hom}_{\Z}\bigl(\pi_0(U\wedge C),\Q/\Z_{(p)}\bigr).
\]

Write \(U\) as a filtered homotopy colimit of finite ordinary
\(p\)-local spectra \(U_i\). Ordinary smash with \(C\) and ordinary
homotopy groups preserve this colimit, so
\[
 \pi_0(U\wedge C)=\varinjlim_i\pi_0(U_i\wedge C).
\]
The functional corresponding to \(f\) vanishes on every term by
the assumed finite precompositions. It therefore vanishes on the
colimit, and \(f=0\). This argument tests by ordinary finite
spectra and uses no interchange of \(K(a)\)-localization with
this filtered colimit.
\end{proof}

\begin{proof}[Proof of Proposition~\ref{wup:upper-null}]
First obtain some finite exponent for the actual local unit,
independently of any finite test. Devinatz--Hopkins prove that
every \(K(h)\)-local spectrum belongs to the thick tensor ideal
generated by \(E_h\) in \(\mathcal C_h\)
\cite[Appendix~A, Proposition~A.3]{dh2004}. This is the local
nilpotence consequence of Hopkins--Ravenel; its descendability
form is \cite[Proposition~10.10]{mathew2016}. The objects of finite
\(E_h\)-exponent form a thick tensor ideal containing \(E_h\):
modules have exponent at most one, tensoring preserves a bound,
and cofiber sequences and retracts preserve finite bounds by
Lemma~\ref{lem:exponent-calculus}. It follows that \(B_h\) has
some finite exponent \(c\). This is one bound on an actual map
before any test is chosen; the nilpotence theorem supplies no
numerical value needed here.

Apply the full local-unit Amitsur construction in \(\mathcal C_h\)
to \(B_h\), and then the exact ordinary functor \(F(V_0,-)\) for a
finite ordinary \(V_0\). Proposition~\ref{wup:cochains} identifies
its second page with
\[
 H^s_{\cts}\bigl(\G_h,\pi_q(E_h\wedge D V_0)\bigr).
\]
Proposition~\ref{prop:cohomology} makes this page zero for
\(s\geq h^2+2=m\). The tested-image assertion of
Proposition~\ref{prop:pages}(ii) applies using the already known
finite exponent \(c\). It concerns the full finite partial
totalizations and their finite fiber cubes, and gives
\[
 [V_0,I_h^{\otimes_h m}]\longrightarrow[V_0,B_h]
       \quad\text{zero for every finite ordinary }V_0.
\]
Thus the actual arrow \(e_m(B_h)\), regarded in ordinary spectra,
vanishes after every finite ordinary precomposition.
Lemma~\ref{wup:phantom} makes it null as an ordinary map and
hence null in the full subcategory \(\mathcal C_h\).

For clarity, all layers needed next occur at finite stages. Set
\(I_h^{\otimes_h0}=B_h\), and for \(0\leq j\leq m\) put
\[
 F_j=\cofib\bigl(
 I_h^{\otimes_h m}\longrightarrow I_h^{\otimes_h(m-j)}
 \bigr).
\]
The consecutive maps form a finite filtration
\(0=F_0\to F_1\to\cdots\to F_m\).
The cofiber of \(I_h^{\otimes_h(i+1)}\to I_h^{\otimes_h i}\)
is \(E_h\otimes_h I_h^{\otimes_h i}\). The cofiber comparison
for consecutive composites therefore identifies
\[
 \cofib(F_j\to F_{j+1})
 \simeq E_h\otimes_h I_h^{\otimes_h(m-j-1)}
       \quad(0\leq j<m).
\]
Finally \(F_m=\cofib(e_m(B_h))\). A nullhomotopy of \(e_m(B_h)\)
makes \(B_h\) a retract of this cofiber. This proves both assertions.
\end{proof}

\section{Field-valued points and coherent trace}\label{sec:coherent-trace}

Continue with $p\geq5$, $t=p-1$, $h=p$, and the pure subgroup $H$
of Theorem~\ref{thm:lower-hm}. To transport the finite upper layers,
we need exponent one for modules over the mixed algebra
\[
 D=E_t\wedge E_h,\qquad D'=L_{K(t)}D.
\]
Here $H$ acts on the first factor. Section~\ref{sec:exponent-transport}
will make each transformed layer a $D'$-module in $\mathcal D_H$.
The \(K(t)\)-acyclic spectra form a smash ideal. Monoidal localization
of \(D\) into \(\mathcal C_t\), followed by the lax monoidal inclusion
into ordinary spectra, therefore gives \(D'=L_{K(t)}D\) a commutative
algebra structure in \(\mathcal D_H\). Its localization map is
equivariant and is a map of these commutative algebras.

\begin{proposition}[The completed mixed trace]\label{wup:completed-trace}
The mixed $H$-algebra $D'$ has an underlying degree-zero class whose
$H$-trace is a unit.
Every \(D'\)-module in \(\mathcal D_H\) therefore has
\(H\)-exponent at most one.
\end{proposition}

A trace unit gives the required coherent retraction by induction and
coinduction. We first prove that implication. We then construct the
trace unit from the two formal-group isomorphisms supplied by the
factors of $D$: the fixed $E_h$ factor forces $H$ to act freely on
field-valued points of
\[
 R=\pi_0D/\m_t\pi_0D.
\]
The ideal is $H$-stable. The mixed cooperation calculation
of Lemma~\ref{wup:mixed-flat} makes $D$ even periodic, with an
orientation and periodic unit from $E_t$.

\begin{lemma}[Coherent trace splitting]\label{wup:trace-split}
Let \(D'\) be a commutative algebra in \(\mathcal D_H\). Suppose an
underlying class \(z\in\pi_0D'\) has unit trace
\[
 v=\sum_{g\in H}g(z)\in(\pi_0D')^\times.
\]
Then every \(D'\)-module \(U\) in \(\mathcal D_H\) has
\(\expn_H(U)\leq1\).
\end{lemma}

\begin{proof}
Since \(H\) is finite, the underlying object \(A_H\wedge U\) is
\(F(H_+,U)\), with diagonal action.
Untwisting the value at \(g\) by the action of \(g^{-1}\) on \(U\)
identifies it with the coinduction of the underlying spectrum:
\[
 \theta:A_H\wedge U\xrightarrow{\ \simeq\ }
       \operatorname{Coind}_{1}^{H}(\operatorname{Res}_{1}^{H}U).
\]
Indeed, on functions the diagonal action is
\((h f)(g)=h f(h^{-1}g)\), and \((\theta f)(g)=g^{-1}f(g)\)
turns it into translation \((h\theta f)(g)=(\theta f)(h^{-1}g)\).
For finite \(H\), the finite coproduct-product equivalence in spectra
identifies induction with coinduction. Induction adjunction takes
the underlying multiplication map \(\mu_z:U\to U\) to a map with
coherent \(H\)-action
\[
 r_z:A_H\wedge U\longrightarrow U.
\]

Let \(j:U\to A_H\wedge U\) be the unit map. Under \(\theta\), its
component at \(g\) is \(g^{-1}:U\to U\). The induction map has
component \(g\mu_z\) at \(g\), so the underlying composite is
\[
 r_zj=\sum_{g\in H}g\mu_zg^{-1}
     =\mu_{\sum_g g(z)}=\mu_v.
\]
The middle equality uses the compatibility of the \(D'\)-module
structure with the \(H\)-action. Since \(v\) is a unit, this is an
equivalence of underlying spectra. Equivalences of \(H\)-diagrams
are detected on underlying objects, so \(r_zj\) is an equivalence
in \(\mathcal D_H\). Hence \((r_zj)^{-1}r_z\) retracts \(j\)
there. Its fiber map is null, which is exponent at most one.
The class \(z\) was only an underlying class; no homotopy-fixed
lift of it is needed.
\end{proof}

\begin{lemma}[Freeness on field-valued points]\label{wup:field-points}
For every unital ring map \(x:R\to\Omega\) to a field, the equality
\(x\circ g=x\) for \(g\in H\) implies \(g=1\).
\end{lemma}

\begin{proof}
We first recall explicitly the naturality of the formal group used
here. For a complex-oriented even periodic commutative ring
spectrum \(B\), an orientation and a periodic unit give a coordinate
in degree zero. The finite projective-space calculation and its
surjective restriction tower give
\[
 B^0(\mathbb{CP}^{\infty})=(\pi_0B)[[z_B]].
\]
Products of projective spaces give the corresponding group law.
A unital ring-spectrum map \(f:B\to C\) carries an orientation
to an orientation and a periodic unit to a unit. Its coordinate
therefore has an invertible linear coefficient in the chosen
coordinate of \(C\). The inverse of this coordinate substitution
gives an isomorphism from the base change of the formal group of
\(B\) to the formal group of \(C\).
The product calculation gives compatibility with group laws,
and this construction respects composites of ring-spectrum maps.
All later base changes are base changes of these algebraic formal
groups along maps of coefficient rings.

Let \(\rho:\pi_0D\to\Omega\) be the composite through \(x\).
Write \(\iota_t:E_t\to D\) and \(\iota_h:E_h\to D\) for the
two factor maps, and let \(\rho_t,\rho_h\) be their coefficient
composites with \(\rho\). The map \(\rho_t\) factors through
\[
 O_t/\m_t=k_t\longrightarrow\Omega.
\]
The last map is injective, since it is a unital map from a field.
After base change along \(\rho_t,\rho_h,\rho\), the factor maps
give formal-group isomorphisms
\[
 \phi_t:\mathcal F_{t,\Omega}\xrightarrow{\ \cong\ }
              \mathcal F_{D,\Omega},\qquad
 \phi_h:\mathcal F_{h,\Omega}\xrightarrow{\ \cong\ }
              \mathcal F_{D,\Omega}.
\]

Let \(g_D\) denote the action of \(g\) on \(D\).
The equality \(x\circ g=x\) implies
\(\rho\circ\pi_0(g_D)=\rho\), so naturality makes \(g_D\) an
automorphism \(u_g\) of the single formal group
\(\mathcal F_{D,\Omega}\). Since \(g_D\iota_h=\iota_h\),
naturality with the second factor gives \(u_g\phi_h=\phi_h\).
Thus \(u_g\) is the identity. But \(g_D\iota_t=\iota_t g\);
naturality with the first factor gives
\[
 u_g\phi_t=\phi_t\,g_\Omega,
\]
where \(g_\Omega\) is the specialization of the prescribed
stabilizer automorphism. Consequently \(g_\Omega\) is the identity.

On the special fiber the pure stabilizer element \(g\) is its
specified automorphism of the Honda formal group over \(k_t\).
This is the special-fiber content of Lubin--Tate deformation
\cite[Proposition~3.3 and paragraph~3.4]{lubin-tate1966}, with its
coherent realization as above. If \(g\ne1\), some coefficient of
the power series \(g(X)-X\) is nonzero in \(k_t\). It remains
nonzero under the injective map \(k_t\to\Omega\), contrary to
\(g_\Omega=1\). Therefore \(g=1\).

This proof used only the coefficient map to \(\Omega\) for algebraic
base change; it did not require a ring-spectrum map to a spectrum
with coefficient field \(\Omega\). The assertion concerns equality
of actual field-valued maps; invariance of their kernels alone
would not give the conclusion needed next.
\end{proof}

\begin{lemma}[Algebraic trace one]\label{wup:algebraic-trace}
Let a finite group \(H\) act on a commutative ring \(R\). Suppose
that for every unital field-valued map \(x:R\to\Omega\),
\(x\circ g=x\) implies \(g=1\). Then
\[
 1\in\operatorname{im}\bigl(\Tr:R\to R^H\bigr),
 \qquad \Tr(r)=\sum_{g\in H}g(r).
\]
\end{lemma}

\begin{proof}
The trace image is an ideal of \(R^H\), since trace is additive
and \(a\Tr(r)=\Tr(ar)\) for \(a\in R^H\). Suppose this ideal
is proper, and choose a maximal ideal \(\mathfrak q\) containing
it. For every \(r\in R\), the monic orbit polynomial
\(\prod_{g\in H}(X-g(r))\) has coefficients in \(R^H\).
Thus \(R\) is integral over \(R^H\). Lying over gives a prime
\(\mathfrak p\) of \(R\) above \(\mathfrak q\). The quotient and
inclusion into its fraction field give a unital map
\(x:R\to\Omega=\operatorname{Frac}(R/\mathfrak p)\) killing
\(\mathfrak q\).

The multiplicative unital maps
\(\chi_g:R\to\Omega\), \(\chi_g(r)=x(g(r))\), are pairwise
distinct by the hypothesis. Such distinct maps are linearly
independent as functions. To prove this, suppose
\(\sum_{i=1}^v a_i\chi_i=0\) is a nonzero relation with the
least possible number of nonzero coefficients. One term is
impossible by evaluation at \(1\). If \(v>1\), choose an element
\(b\in R\) distinguishing \(\chi_1\) from \(\chi_v\).
Subtract \(\chi_v(b)\) times the relation at \(r\) from the
relation at \(br\). The result is
\[
 \sum_{i<v}a_i\bigl(\chi_i(b)-\chi_v(b)\bigr)\chi_i(r)=0.
\]
Its coefficient at \(i=1\) is nonzero in the field, giving a
shorter nonzero relation, a contradiction.

In particular the function \(\sum_{g\in H}\chi_g\) is not zero:
each of its coefficients is \(1\) in the field, even if the
characteristic divides \(|H|\). On the other hand it equals
\(x\circ\Tr\), which is zero since \(x\) kills the trace ideal.
This contradiction proves the claim. If \(R\) is the zero ring,
the conclusion holds directly with \(1=0\).
\end{proof}

\begin{proof}[Proof of Proposition~\ref{wup:completed-trace}]
Lemmas~\ref{wup:field-points} and \ref{wup:algebraic-trace}
give \(\bar z\in R\) with \(\Tr(\bar z)=1\). Lift it to
\(z_0\in\pi_0D\), and put \(v=\sum_g g(z_0)\). Then
\[
 v-1\in\m_t\pi_0D.
\]
The mixed module \(\pi_*D\) is flat over \((E_t)_*\) by
Lemma~\ref{wup:mixed-flat}. Lemma~\ref{wup:completion-reduction}
therefore identifies the homotopy of its actual localization
\(D\to D'\) with the ordinary \(\m_t\)-adic completion of
\(\pi_*D\).

At the finite adic stage \(\pi_0D/\m_t^j\pi_0D\), the element
\[
 w_j=\sum_{k=0}^{j-1}(1-v)^k
\]
satisfies \(v w_j=1-(1-v)^j=1\). These inverses are compatible
under reduction. Here \(j\) is a finite truncation level; the
quotient rings themselves are not required to be finite sets.
The comparison with completion is induced by the actual ring map
\(D\to D'\). Multiplication by the image of \(v\) is
\(O_t\)-linear and adically continuous and agrees with
multiplication by \(v\) on the dense ordinary image. It hence
agrees with the levelwise multiplication on the completion.
The compatible \(w_j\) give an inverse to the image of \(v\) in
\(\pi_0D'\).

Localization is equivariant, so that image of \(v\) is the
\(H\)-trace of the image of \(z_0\). Lemma~\ref{wup:trace-split}
now applies. If the completed ring is zero, its unit is zero and
every unital module is contractible, giving the same conclusion.
\end{proof}

\section{Transport of the finite layers and the canonical unit}
\label{sec:exponent-transport}

The upper null map supplies a retract built from $m$ local $E_h$-module
layers. We now transport those individual layers by the exact functor
on ordinary \(p\)-local spectra
\[
 T:\Sp\longrightarrow\mathcal D_H,\qquad
 T(V)=L_{K(t)}(E_t\wedge V),
\]
with \(H\) acting through the coherent action on \(E_t\).
Applying \(T\) to a local layer means applying it to the underlying
ordinary spectrum of that object of \(\mathcal C_h\). The mixed trace
will make each resulting layer have $H$-exponent at most one.

\begin{proposition}[The cross-height upper exponent]\label{wup:transfer}
Let \(p\geq5\) be prime, with \(t=p-1\) and \(h=p\). One has
\[
 \expn_H\bigl(T(B_h)\bigr)\leq m=p^2+2.
\]
\end{proposition}

\begin{proof}
The inclusion \(\mathcal C_h\to\Sp\) is exact and lax symmetric
monoidal. Since \(E_h\) is local, the localization unit identifies
the included local algebra \(E_h\) with its given ordinary algebra.
Each local layer
\(E_h\otimes_h I_h^{\otimes_h i}\) in
Proposition~\ref{wup:upper-null} consequently has an ordinary
\(E_h\)-module structure. Ordinary smash with \(E_t\) makes it
a \(D=E_t\wedge E_h\)-module with compatible \(H\)-action.

The \(K(t)\)-acyclic spectra form a smash ideal. Thus the
localization into \(\mathcal C_t\) is symmetric monoidal for its
localized tensor, and the inclusion back into ordinary spectra
is lax symmetric monoidal. Applying localization to the algebra
and this module gives, after that inclusion, a \(D'\)-module
structure on the transformed layer in \(\mathcal D_H\).
Proposition~\ref{wup:completed-trace} gives every such layer
\(H\)-exponent at most one.

The functor \(T\) preserves cofiber sequences and retracts.
Apply it to the finite filtration and retraction of
Proposition~\ref{wup:upper-null}, and add the \(m\) layer bounds
using Lemma~\ref{lem:exponent-calculus}. This yields the stated
bound. The construction used ordinary module structures on the
individual layers; it required no preservation by \(T\) of
the \(\mathcal C_h\) tensor product.
\end{proof}

\begin{corollary}[The exponent contradiction for the canonical unit]
\label{wup:contradiction}
Let \(p\geq5\) be prime, put \(t=p-1\) and \(h=p\), and let
\(X=S^\wedge_p\) be the derived \(p\)-completed sphere. The canonical map
\[
 i_{h,X}=L_t(\eta_X^{(h)}):
 L_tX\longrightarrow L_tL_{K(h)}X
\]
has no homotopy retraction in ordinary spectra.
\end{corollary}

\begin{proof}
The functor \(T\) inverts \(K(t)\)-equivalences: the cofiber of
such a map is \(K(t)\)-acyclic, and remains so after ordinary
smash with \(E_t\). Every \(L_t\)-localization unit is a
\(K(t)\)-equivalence. Lemma~\ref{lem:completion} says that the
actual completion \(c:S\to X\) is a \(K(a)\)-equivalence for
each positive \(a\), in particular for \(a=t,h\).
Write $\lambda_Z:Z\to L_tZ$ for the lower localization unit.
Consider the square of actual localization and completion maps
\[
\begin{CD}
 S @>{\eta_S^{(h)}}>> B_h\\
 @V{\lambda_X\,c}VV
 @VV{\lambda_{L_{K(h)}X}\,L_{K(h)}(c)}V\\
 L_tX @>{i_{h,X}}>> L_tL_{K(h)}X .
\end{CD}
\]
It commutes by naturality of the two localization units.
Both vertical arrows become equivalences under \(T\): for the
left one use the \(K(t)\)-equivalences just stated, and for the
right one use that \(L_{K(h)}(c)\) is an equivalence followed
by a \(K(t)\)-equivalence. Moreover \(T(S)\simeq E_t\), with
coherent \(H\)-action, since \(E_t\) is \(K(t)\)-local.

A homotopy retraction of \(i_{h,X}\) would therefore, after
applying \(T\) and these identifications, exhibit \(E_t\) as a
retract of \(T(B_h)\) in \(\mathcal D_H\).
Proposition~\ref{wup:transfer} and preservation of exponent
bounds under retracts would give
\(\expn_H(E_t)\leq m\). The finite-page lower result
Proposition~\ref{prop:lower-exponent} gives the strict inequality
\(\expn_H(E_t)>m\), a contradiction.
\end{proof}

\begin{proof}[Proof of Theorem~\ref{thm:exponents}]
Proposition~\ref{prop:lower-exponent} gives the strict lower bound and
Proposition~\ref{wup:transfer} gives the upper bound, both for
$m=p^2+2$ in the single category $\mathcal D_H$. The preceding corollary
proves that a retraction of the canonical map would turn the lower object
into an $H$-equivariant retract of the upper object and contradict those
bounds. This proves every assertion of the theorem. The lower bound uses
finite-page survival only; the upper bound uses the actual null map and
its finite layers.
\end{proof}

\appendix
\section{Ordinary products in homotopy fixed points}\label{sec:ordinary-products}
We prove Lemma~\ref{lem:lower-bar-product} by constructing relative
products on the ordinary bar tower. The construction supplies cup
products and differential derivations independently of convergence;
separated convergence identifies the final product with the product
on actual homotopy groups.

\begin{proof}[Proof of Lemma~\ref{lem:lower-bar-product}]
Use the homogeneous free bar construction $EQ_\bullet$, with
$EQ_q=Q^{q+1}$, and let $EQ^{(s)}$ be its $s$-skeleton.  The full
totalization tower of the group-action resolution is
\[
 V_s=F_Q((EQ^{(s)})_+,R),\qquad
 \operatorname{holim}_sV_s=F_Q(EQ_+,R)=R^{hQ}.
\]
Here $F_Q$ denotes the derived spectrum of equivariant maps.  The
skeletal filtration gives the usual homogeneous cochains and hence the
displayed $E_2$ page.

For the prismatic geometry below, compare Basterra--Mandell
\cite[proof of Theorem~7.1, pp.~23--26]{basterra-mandell2013}.
Their bar diagonal for a (partial) non-unital $E_1$ algebra over a base
ring spectrum $H$ uses the
same half-cut regions and inverse barycentric coordinates. Here we
construct the equivariant relative maps for $EQ$ and identify the
resulting $E_\infty$ product with the associated-graded product on actual
homotopy groups under the lemma's stated filtration hypothesis.
For barycentric coordinates $t=(t_0,\ldots,t_q)\in\Delta^q$, set
$c_{-1}=0$ and $c_i=t_0+\cdots+t_i$.  Define two barycentric vectors by
\begin{equation}\label{eq:lower-half-cut}
 \begin{split}
 \ell_i(t)&=2\bigl(\min(c_i,\tfrac12)-\min(c_{i-1},\tfrac12)\bigr),\\
 \rho_i(t)&=2\bigl(\max(c_i,\tfrac12)-\max(c_{i-1},\tfrac12)\bigr).
 \end{split}
\end{equation}
They are nonnegative and each has sum one.  Equivalently, their $i$th
coordinates are twice the lengths of the intersections of
$[c_{i-1},c_i]$ with $[0,\tfrac12]$ and $[\tfrac12,1]$.  The formulas
are continuous, including when a cumulative sum equals $\tfrac12$.

For an order-preserving map $\theta:[q]\to[q']$, let
$(\theta_*t)_j=\sum_{\theta(i)=j}t_i$.  Each interval belonging to
$\theta_*t$ is the union of the consecutive intervals belonging to the
corresponding fiber of $\theta$; an empty fiber gives an interval of
length zero.  Additivity of intersection lengths gives
\[
 \ell(\theta_*t)=\theta_*\ell(t),\qquad
 \rho(\theta_*t)=\theta_*\rho(t).
\]
The realization relation is
$[\theta^*\sigma,t]=[\sigma,\theta_*t]$ for $\sigma\in EQ_{q'}$.
Thus for every bar simplex $\sigma\in EQ_q$ the formula
\[
 \delta[\sigma,t]=([\sigma,\ell(t)],[\sigma,\rho(t)])
\]
is compatible with all faces and degeneracies and defines a continuous
map $\delta:EQ\to EQ\times EQ$.  It is $Q$-equivariant because the
action changes $\sigma$ and leaves the barycentric coordinates fixed.
Moreover
\[
 [\sigma,t]\longmapsto
 \bigl([\sigma,(1-\tau)t+\tau\ell(t)],
       [\sigma,(1-\tau)t+\tau\rho(t)]\bigr),
 \qquad 0\leq\tau\leq1,
\]
is compatible with the same identifications and is equivariant.  It is
a homotopy from the ordinary diagonal to $\delta$.

Choose $k$ with $c_{k-1}\leq\tfrac12\leq c_k$.  The support of
$\ell(t)$ is contained in the front face $[0,\ldots,k]$, and the support
of $\rho(t)$ is contained in the back face $[k,\ldots,q]$.  These faces
have dimensions $k$ and $q-k$.  Degenerate faces can only lower those
dimensions.  Consequently
\[
 \delta(EQ^{(s)})\subset
 \bigcup_{i+j\leq s}EQ^{(i)}\times EQ^{(j)}.
\]
This also identifies the cellular chain map.  On the region
$c_{k-1}\leq\tfrac12\leq c_k$, the map
$t\mapsto(\ell(t),\rho(t))$ is a homeomorphism onto
$\Delta^k\times\Delta^{q-k}$, with inverse
\[
 t_i=\ell_i/2\ (i<k),\qquad
 t_k=(\ell_k+\rho_k)/2,\qquad
 t_i=\rho_i/2\ (i>k).
\]
In the oriented coordinates $t_1,\ldots,t_q$ on the source and
$\ell_1,\ldots,\ell_k,\rho_{k+1},\ldots,\rho_q$ on the product, its
Jacobian determinant is $2^q>0$.  The regions meet on their boundaries.
The image of the cellular fundamental chain is therefore
\[
 [\sigma]\longmapsto
 \sum_{k=0}^q
 [\sigma|_{[0,\ldots,k]}]\otimes
 [\sigma|_{[k,\ldots,q]}],
\]
with degenerate terms omitted in normalized chains.  This is the
Alexander--Whitney diagonal, now obtained from the explicit map.

Write $EQ/EQ^{(s)}$ for the based cofiber of
$(EQ^{(s)})_+\to EQ_+$, and set $EQ^{(-1)}=\varnothing$.
The skeletal inclusions are $Q$-CW cofibrations, so these quotients
represent their homotopy cofibers.  For $a,b\geq0$, the skeletal
containment above makes the composite of $\delta$ with the quotient to
$(EQ/EQ^{(a-1)})\wedge(EQ/EQ^{(b-1)})$ constant on
$EQ^{(a+b-1)}$: if $i+j<a+b$, then $i<a$ or $j<b$.  We get a
$Q$-equivariant relative map
\[
 EQ/EQ^{(a+b-1)}\longrightarrow
 (EQ/EQ^{(a-1)})\wedge(EQ/EQ^{(b-1)}).
\]
Multiplication on $R$ and precomposition give
\[
 \begin{split}
 F_Q(EQ/EQ^{(a-1)},R)\wedge F_Q(EQ/EQ^{(b-1)},R)
 \longrightarrow F_Q(EQ/EQ^{(a+b-1)},R).
 \end{split}
\]
These maps are compatible as $a$ and $b$ vary, since they are induced
by the same $\delta$ and the nested skeletal quotients.

For completeness, the associativity needed for powers also holds within
this relative filtration.  Partition $[0,1]$ into any $m$ successive
intervals of positive lengths.  For the $j$th output vector, divide
the lengths of its intersections with $[c_{i-1},c_i]$ by the length of
the $j$th interval.  The preceding naturality and equivariance proof
applies verbatim.  If $k_j$ is an index whose coordinate interval
contains the $j$th cut, the supports lie in the successive faces
$[0,\ldots,k_1]$, $[k_1,\ldots,k_2]$, \ldots,
$[k_{m-1},\ldots,q]$.  Their dimensions sum to $q$; boundary cases
can only reduce the sum.  Thus the same construction gives $m$-fold
relative maps.  Composing cut maps is exactly refinement of the
partition.  The positive interval lengths form a convex simplex, so
the partitions for two parenthesizations are joined by varying those
lengths.  This gives a homotopy through maps with the same relative
skeletal property.  Thus the relative products are associative up to
filtered homotopy, as required for their iterated products.

Put $I_s=F_Q(EQ/EQ^{(s-1)},R)$ for $s\geq0$.  The quotient cofiber
sequence for two successive skeleta gives
\[
 \operatorname{cofib}(I_{s+1}\to I_s)
 \simeq F_Q(EQ^{(s)}/EQ^{(s-1)},R)
 \simeq \operatorname{fib}(V_s\to V_{s-1}).
\]
The fiber sequences $I_s\to R^{hQ}\to V_{s-1}$ form a diagram with
the identity on $R^{hQ}$ and the tower transition on the right.
The stable fiber diagram identifies the connecting maps of the
displayed cofiber sequences with those of the $V_s$ tower, with the
usual rotation signs.  Thus the paired relative spectral sequence is
the ordinary tower spectral sequence with the same filtration index.
The relative products pair these sequences and their exact couples.
On cellular cochains the displayed chain formula is the ordinary cup
product; in inhomogeneous coordinates its value is
\[
 (v\smile w)(g_1,\ldots,g_{a+b})
 =v(g_1,\ldots,g_a)\,(g_1\cdots g_a)
       w(g_{a+1},\ldots,g_{a+b}),
\]
with the usual graded signs.  The cellular boundary formula for a
product, together with the internal Koszul sign, gives the derivation
rule in the paired exact couples and hence on every spectral-sequence
page.  This $E_2$ product is the ordinary cup product used by the
ordinary-to-Tate comparison in Section~\ref{sec:lower-exponent}.
Since each subsequent product
is induced on homology, the products agree on every later page as well.

The fiber sequence for $I_a$ identifies its image in $\pi_*R^{hQ}$ with
$\ker(\pi_*R^{hQ}\to\pi_*V_{a-1})$. After
forgetting the relative conditions, the product just constructed is
homotopic to the actual product on $F_Q(EQ_+,R)$ by the equivariant
homotopy from $\delta$ to the diagonal.  Thus these relative products
multiply the actual filtration groups.  Passing to successive relative
quotients is exactly the product in the paired spectral sequence.
When the stated convergence identifies those quotients with
$E_\infty$, this is the associated-graded homotopy product as stated.
\end{proof}

\clearpage
\phantomsection
\addcontentsline{toc}{section}{References}
\bibliographystyle{plain}
\bibliography{references}
\end{document}
